% =====================================================================================================
%  پایان‌نامه: بررسی، تحلیل و بهینه‌سازی برهم‌کنش گسیلنده‌ی کوانتومی منفرد با نانوآنتن پلاسمونی طلا
%  کامپایلر Overleaf باید XeLaTeX باشد:  Menu -> Settings -> Compiler -> XeLaTeX
%  تصاویر (اختیاری): purcell_nm.png, box_x1_nm.png, box_x2_nm.png, box_y1_nm.png, box_y2_nm.png, box_z1_nm.png, box_z2_T.png
%  اگر فایل تصویری بارگذاری نشده باشد، به جای آن یک کادر خالی چاپ می‌شود و خطایی رخ نمی‌دهد.
% =====================================================================================================
\documentclass[12pt,a4paper]{report}

\usepackage{amsmath,amssymb}
\usepackage{graphicx}
\usepackage[a4paper,top=2.8cm,bottom=2.6cm,right=2.6cm,left=2.4cm,headheight=16pt]{geometry}
\usepackage{setspace}
\usepackage{array}
\usepackage{booktabs}
\usepackage{enumitem}
\usepackage{fancyhdr}
\usepackage[colorlinks=true,linkcolor=black,citecolor=blue,urlcolor=blue]{hyperref}
\usepackage{xepersian}   % باید آخرین بسته باشد

% ---------------------------------------- قلم یاس ----------------------------------------------------
% به ترتیب امتحان می‌شود: فایل قلم بارگذاری‌شده در پروژه، سپس قلم‌های نصب‌شده با نام‌های رایج «یاس».
% Scale=1.12 متن را حدود ۱۲ درصد بزرگ‌تر از نسخه‌ی قبلی می‌کند.
% اگر در پروژه‌ی قبلی‌تان خطی مثل \settextfont{...} داشتید، می‌توانید کل این بلوک را با همان خط
% (با افزودن Scale=1.12) جایگزین کنید.
\IfFontExistsTF{BYas.ttf}{%
  \IfFontExistsTF{BYasBd.ttf}%
    {\settextfont[Scale=1.12,BoldFont=BYasBd.ttf]{BYas.ttf}}%
    {\IfFontExistsTF{BYasBold.ttf}%
      {\settextfont[Scale=1.12,BoldFont=BYasBold.ttf]{BYas.ttf}}%
      {\settextfont[Scale=1.12]{BYas.ttf}}}}%
{\IfFontExistsTF{BYas}{\settextfont[Scale=1.12]{BYas}}%
{\IfFontExistsTF{B Yas}{\settextfont[Scale=1.12]{B Yas}}%
{\IfFontExistsTF{XB Yas}{\settextfont[Scale=1.12]{XB Yas}}%
{\IfFontExistsTF{Yas}{\settextfont[Scale=1.12]{Yas}}%
{\IfFontExistsTF{XB Niloofar}{\settextfont[Scale=1.12]{XB Niloofar}}%
{\IfFontExistsTF{Vazirmatn}{\settextfont[Scale=1.05]{Vazirmatn}}%
{\IfFontExistsTF{Nazli}{\settextfont[Scale=1.12]{Nazli}}%
{\settextfont[Scale=1.05]{Amiri}}}}}}}}}

% ---------------------------------------- فاصله‌ی خطوط --------------------------------------------
\setstretch{1.55}
\setlength{\parskip}{3pt}
\setlist{itemsep=4pt,topsep=4pt}
\setlength{\emergencystretch}{3em}

% ---------------------------------------- سرصفحه ---------------------------------------------------
\pagestyle{fancy}
\fancyhf{}
\fancyhead[R]{\small\leftmark}
\fancyfoot[C]{\thepage}
\renewcommand{\headrulewidth}{0.5pt}
\fancypagestyle{plain}{\fancyhf{}\fancyfoot[C]{\thepage}\renewcommand{\headrulewidth}{0pt}}

% ---------------------------------------- فرمان‌های کمکی ------------------------------------------
\newcommand{\sub}[1]{\mathrm{#1}}               % زیرنویس لاتین در فرمول‌ها (مثل rad، tot)
\newcommand{\Imag}{\operatorname{Im}}
\newcommand{\Real}{\operatorname{Re}}
% محل تصویر: اگر فایل وجود داشته باشد درج می‌شود، وگرنه کادر خالی با نام فایل
\newcommand{\figbox}[2]{%
  \IfFileExists{#1}{\includegraphics[width=#2]{#1}}{%
    \fbox{\parbox[c][\dimexpr#2*42/100\relax][c]{\dimexpr#2-2\fboxsep-2\fboxrule\relax}{\centering\textbf{محل قرارگیری نمودار لومریکال}\\[4pt]
    فایل تصویر را با نام \lr{\texttt{\detokenize{#1}}} بارگذاری کنید.}}}}

\begin{document}

% ======================================== صفحه‌ی عنوان ============================================
\begin{titlepage}
\begin{center}
\vspace*{1.5cm}
{\Large\textbf{دانشگاه اصفهان}}\\[6pt]
{\large دانشکده‌ی فیزیک}\\[4pt]
\rule{\textwidth}{0.6pt}\\[10pt]
{\LARGE\textbf{بررسی، تحلیل و بهینه‌سازی برهم‌کنش گسیلنده‌ی کوانتومی منفرد با نانوآنتن پلاسمونی طلا
(نانومیله‌ی کلاهک‌دار) به روش \lr{3D-FDTD}}}\\[14pt]
{\large\textbf{تفکیک کانال‌های تابشی و غیرتابشی و تعیین فاصله‌ی بهینه برای چشمه‌های تک‌فوتونی}}\\[10pt]
\rule{\textwidth}{0.6pt}\\[2.5cm]
\begin{minipage}[t]{0.48\textwidth}
\textbf{پژوهشگر:}\\
سید محمد پارسا مولایی طبری\\[12pt]
\textbf{شماره‌ی دانشجویی:}\\
۴۰۲۴۲۱۳۰۴۶
\end{minipage}\hfill
\begin{minipage}[t]{0.42\textwidth}
\textbf{استاد راهنما:}\\
دکتر مالک باقری هارونی
\end{minipage}
\end{center}
\end{titlepage}

% ======================================== چکیده =====================================================
\chapter*{چکیده}
\addcontentsline{toc}{chapter}{چکیده}
\markboth{چکیده}{چکیده}
کنترل نرخ گسیل خودبه‌خودیِ گسیلنده‌های کوانتومی به کمک محیط‌های ساختاریافته، نقشی اساسی در توسعه‌ی چشمه‌های تک‌فوتونیِ کارآمد برای فناوری‌های اطلاعات کوانتومی، رمزنگاری کوانتومی و مخابرات نوری دارد. نانوآنتن‌های پلاسمونی به دلیل توانایی‌شان در متمرکز کردن میدان‌های نوری در مقیاس‌های کوچک‌تر از حد پراش، گزینه‌هایی مناسب برای تقویت جفت‌شدگی نور و ماده (اثر پورسل) به شمار می‌روند. با این حال، نزدیکی گسیلنده به فلزات واقعی همواره با یک رقابت فیزیکی همراه است: رقابت میان تقویت نرخ تابشی و اتلاف توان به صورت گرما (خاموشی غیرتابشی) در فواصل نانومتری.

در این پژوهش، برهم‌کنش یک گسیلنده‌ی کوانتومی (مدل‌شده به صورت نوسانگر دوقطبی کلاسیک) با یک نانومیله‌ی طلای کلاهک‌دار (به طول ۶۰ و قطر ۲۰ نانومتر) با روش تفاضل محدود در حوزه‌ی زمان (\lr{3D-FDTD})، که معادلات ماکسول را بدون تقریب و به صورت عددی حل می‌کند، به صورت سه‌بعدی شبیه‌سازی شد. هدف اصلی، تفکیک سهم کانال‌های تابشی و غیرتابشی و یافتن فاصله‌ی بهینه‌ی گسیلنده تا سطح فلز (گاف) برای بیشینه‌سازی بازده نانوآنتن بوده است. برای سنجش نقش هندسه، نتایج نانومیله در شرایط یکسان (حجم برابر، جنس، محیط، جهت دوقطبی و فاصله‌ی یکسان) با یک نانوکره‌ی طلا، که مرجع کلاسیک این حوزه است، مقایسه شد.

نتایج تحلیل طیفی نشان داد که نانومیله‌ی طلا ناهمسانگردی جهتی شدیدی دارد و تشدید پلاسمون سطحی طولی آن در طول‌موج حدود ۶۰۷ نانومتر رخ می‌دهد. در فواصل بسیار نزدیک (۳ تا ۵ نانومتر)، اتلاف اهمی فلز غالب است و بازده کوانتومی ظاهری حداکثر به ۱۹ درصد می‌رسد. با افزایش فاصله به فواصل میانی (حدود ۱۰ تا ۲۰ نانومتر)، سامانه به یک مصالحه‌ی مطلوب می‌رسد: نرخ تابشی همچنان ده‌ها برابر فضای آزاد تقویت می‌شود (حدود ۶۷ برابر در فاصله‌ی ۲۰ نانومتر) و بازده کوانتومی ظاهری تا حدود ۷۴ درصد افزایش می‌یابد. مقایسه با نانوکره‌ی هم‌حجم نشان داد که در فاصله‌ها و طول‌موج یکسان، تقویت تابشی نانومیله حدود ۳۵ تا ۶۰ برابر نانوکره است و تشدید نانومیله، برخلاف تشدید ثابت و پراتلاف نانوکره در حدود ۵۳۰ نانومتر، با نسبت طول به عرض قابل تنظیم است. این یافته‌ها معیارهای طراحی و فاصله‌گذاری دقیق را برای ساخت چشمه‌های فوتونیِ فوق‌سریع مبتنی بر پلاسمونیک ارائه می‌دهند.

\vspace{6pt}
\noindent\textbf{واژگان کلیدی:} نانوآنتن پلاسمونی، نانومیله‌ی طلا، نانوکره‌ی طلا، اثر پورسل، چگالی حالت‌های موضعی (\lr{LDOS})، خاموشی غیرتابشی، روش تفاضل محدود در حوزه‌ی زمان (\lr{3D-FDTD})، بازده کوانتومی، چشمه‌ی تک‌فوتونی.

\tableofcontents

% ======================================== فصل ۱ ======================================================
\chapter{مقدمه و کلیات پژوهش}

\section{مقدمه}
دستکاری و کنترل نرخ گسیل خودبه‌خودی اتم‌ها، مولکول‌های فلورسنت و نقاط کوانتومی از طریق اصلاح چگالی موضعی حالت‌های فوتونی (\lr{LDOS}) یکی از موضوعات محوری و بنیادین در الکترودینامیک کوانتومی درون‌کاواکی (\lr{CQED}) و نانوفوتونیک مدرن به شمار می‌رود \cite{novotny,maier}. در اوایل قرن بیستم و در چارچوب مکانیک کوانتومی اولیه، تصور بر این بود که گسیل خودبه‌خودی یک ویژگی ذاتی و تغییرناپذیر اتم است که تنها به ساختار ترازهای انرژی درونی آن بستگی دارد و محیط اطراف هیچ تأثیری بر دینامیک فروافت سامانه ندارد \cite{novotny}. با این حال، در سال ۱۹۴۶، ادوارد میلز پورسل (\lr{E. M. Purcell}) در مقاله‌ای پیش‌بینی کرد که نرخ گسیل خودبه‌خودی یک سامانه‌ی کوانتومی می‌تواند با قرارگیری در داخل یا مجاورت یک محیط ساختاریافته (مانند یک تشدیدگر الکتریکی یا کاواک رادیویی) به شدت تغییر یابد؛ پدیده‌ای که امروزه با عنوان «اثر پورسل» (\lr{Purcell Effect}) شناخته می‌شود \cite{purcell,novotny}.

برای دستیابی به اثر پورسل قوی و تغییر ملموس در نرخ‌های گذار، تشدیدگر پیرامون گسیلنده باید دو ویژگی کلیدی داشته باشد: نخست، ضریب کیفیت (\lr{Q-factor}) بالا که نشان‌دهنده‌ی توانایی کاواک در به دام انداختن فوتون‌ها برای مدت زمان طولانی است؛ و دوم، حجم مدی (\lr{Mode volume}) بسیار کوچک که بیانگر میزان جای‌گزیدگی فضایی میدان الکترومغناطیسی در اطراف گسیلنده است \cite{novotny,maier}. در حالی که کاواک‌های دی‌الکتریک سنتی، میکروحلقه‌ها (\lr{Micro-ring resonators}) و بلورهای فوتونی ضریب کیفیت بسیار بالایی دارند و اتلاف نوری آن‌ها ناچیز است، به دلیل محدودیت فیزیکی ناشی از «حد پراش نور» (\lr{Diffraction Limit})، حجم مدی آن‌ها نسبتاً بزرگ است \cite{maier,novotny}. بر اساس معیار رایلی، در محیط‌های دی‌الکتریک شفاف، میدان الکترومغناطیسی را نمی‌توان در ابعادی کوچک‌تر از مرتبه‌ی مکعب طول موج، یعنی $(\lambda/2n)^3$، متمرکز کرد \cite{novotny}. این محدودیت سبب می‌شود جفت‌شدگی میان میدان و گسیلنده‌های نانومتری در کاواک‌های دی‌الکتریک از حد مشخصی فراتر نرود \cite{novotny,maier}.

در این میان، نانوساختارهای فلزی (پلاسمونی) راهکاری کارآمد برای عبور از حد پراش ارائه می‌دهند \cite{maier,novotny}. این نانوساختارها با تکیه بر پدیده‌ی «تشدید پلاسمون‌های سطحی موضعی» (\lr{LSPR})، نور تابشی را با نوسانات جمعی الکترون‌های آزاد فلز (گاز الکترونی فرمی) جفت می‌کنند و میدان‌های الکترومغناطیسی را در مقیاس‌هایی بسیار کوچک‌تر از طول موج (در ابعاد چند نانومتر) جای‌گزیده می‌کنند \cite{maier}. این ویژگی، حجم مدی بسیار کوچکی (بسیار کمتر از $\lambda^3$) فراهم می‌آورد و نانوذرات فلزی را به گزینه‌هایی مناسب برای ساخت «نانوآنتن‌های نوری» تبدیل کرده است \cite{novotny,bharadwaj}. نانوآنتن‌ها به عنوان مبدل‌هایی عمل می‌کنند که انرژی موضعی و به شدت متمرکز یک دوقطبی کوانتومی را دریافت کرده و آن را با بازده بالا به امواج الکترومغناطیسیِ در حال انتشار در فضای آزاد تبدیل می‌کنند \cite{novotny,maier}.

\section{بیان مسئله}
در میان هندسه‌های مختلف نانوذرات فلزی، نانومیله‌های طلا (\lr{Gold Nanorods}) به دلیل ویژگی‌های هندسی و اپتیکی متمایز خود، جایگاه ویژه‌ای در طراحی نانوآنتن‌های پیشرفته دارند \cite{maier,giannini}. برخلاف نانوذرات کروی که به دلیل تقارن کامل، تنها یک بسامد تشدید دوقطبی یگانه و ثابت دارند، نانومیله‌ها با شکستن تقارن کروی و برخورداری از یک محور ممتاز طولی، دو مد تشدیدی کاملاً مجزا (مد طولی و مد عرضی) دارند \cite{maier,novotny}. تشدید نانوکره‌ی طلا در محیط هوا یا خلأ در حدود ۵۲۰ تا ۵۳۰ نانومتر رخ می‌دهد و تقریباً بر لبه‌ی گذارهای بین‌نواری طلا منطبق است؛ به همین دلیل پهن و پراتلاف است و تنها با تغییر اندازه یا محیط، به مقدار اندکی جابه‌جا می‌شود \cite{bohren,kreibig,link}. در مقابل، شکستن تقارن کروی در نانومیله‌ها امکان تنظیم دقیق طول‌موج تشدید طولی را در بازه‌ی وسیع مرئی و فروسرخ نزدیک (\lr{Near-Infrared}) فراهم می‌آورد؛ این تنظیم عمدتاً با تغییر نسبت طول به عرض (\lr{Aspect Ratio}) در فرایند سنتز شیمیایی انجام می‌شود \cite{maier,link}. به همین دلیل، نانوکره‌ی طلا در این پژوهش به عنوان «مرجع» یا گروه کنترل به کار رفته است تا سهم هندسه‌ی میله‌ای در رفتار نانوآنتن به طور جداگانه سنجیده شود.

علاوه بر قابلیت تنظیم‌پذیری طیفی، به دلیل انحنای شدید در قطب‌ها یا نوک‌های نانومیله (کلاهک‌های نیم‌کروی)، خطوط میدان جابه‌جایی الکتریکی در این نواحی به شدت متراکم می‌شوند \cite{maier,novotny}. این تراکم شدید بار سطحی در نقاط با شعاع انحنای کوچک، پدیده‌ای موسوم به «اثر برق‌گیر» (\lr{Lightning-rod effect}) را ایجاد می‌کند \cite{novotny,bharadwaj}. پیامد مستقیم این پدیده، تشکیل «نقاط داغ» (\lr{Hotspots}) پلاسمونی است؛ نواحی‌ای که شدت میدان نزدیک (\lr{Near-field}) در آن‌ها می‌تواند صدها تا هزاران برابر میدان تابشی فرودی باشد \cite{novotny,maier}. در نانوکره نیز میدان نزدیک در دو قطبِ هم‌راستا با قطبش بیشینه می‌شود، اما چون انحنای سطح یکنواخت و شعاع آن نسبتاً بزرگ است، تمرکز اضافی ناشی از نوک‌های تیز وجود ندارد و تقویت میدان به مراتب کمتر است \cite{bohren,maier}.

با این حال، به کارگیری این ساختارها برای کنترل گسیل خودبه‌خودی با یک چالش اساسی و فیزیکی همراه است \cite{novotny,anger}. قرارگیری گسیلنده‌ی کوانتومی در فاصله‌ی بسیار نزدیک به سطح فلز (به منظور قرار گرفتن در کانون نقاط داغ و بهره‌مندی از بیشترین شدت میدان و بالاترین فاکتور پورسل)، پدیده‌ای پیچیده به نام «رقابت میان تقویت تابشی و اتلاف اهمی» را به دنبال دارد \cite{novotny,maier}. در فواصل نانومتری (کمتر از ۱۰ نانومتر)، جفت‌شدگی میدان نزدیک گسیلنده با سطح فلز به قدری شدید است که بخش عمده‌ای از انرژی سامانه، به جای تابش به شکل فوتون‌های مفید، جذب ساختار فلزی می‌شود \cite{novotny,anger}.

انرژی منتقل‌شده به فلز در نهایت از دو سازوکار جذب به گرما تبدیل می‌شود \cite{maier,jc}: نخست، جذب درون‌نواری (\lr{Intraband}) ناشی از برخورد الکترون‌های آزاد نوار رسانش (همان میرایی مدل درود)؛ و دوم، گذارهای بین‌نواری (\lr{Interband}) الکترون‌ها از نوار $5d$ به نوار رسانش $6sp$ در طلا. در فواصل بسیار کوتاه، بخش بزرگی از این انرژی ابتدا به مدهای چندقطبی مرتبه‌بالا و غیرتابشی سطح فلز منتقل می‌شود که توانایی انتشار به فضای دوردست را ندارند و سپس از همین دو مسیر مستهلک می‌شوند \cite{novotny,ruppin}. این اتلاف اهمی شدید، منجر به بروز پدیده‌ی افت ناگهانی بازده یا همان «خاموشی غیرتابشی» (\lr{Quenching}) می‌گردد \cite{novotny,anger}.

مسئله‌ی اصلی این پژوهش آن است که با نزدیک شدن گسیلنده به سطح نانومیله‌ی طلا، اگرچه نرخ کل فروافت (فاکتور پورسل) به صورت توانی (عمدتاً متناسب با $1/d^3$) و تا مرتبه‌ی هزاران برابر رشد می‌کند، اما به دلیل غلبه‌ی نرخ فروافت غیرتابشی، بازده تابشی سامانه به شدت افت می‌کند \cite{novotny,koenderink}. بنابراین، یافتن مصالحه‌ی فیزیکی میان این دو پدیده‌ی متضاد، تفکیک دقیق سهم کانال‌های تابشی و غیرتابشی با استفاده از روش‌های عددی حل کامل معادلات ماکسول، و در نهایت «یافتن فاصله‌ی (گاف) بهینه» برای غلبه بر اتلاف اهمی و بیشینه‌سازی خروجی تابشی، مسئله‌ی اصلی این پژوهش است \cite{novotny,maier}. پاسخ به این پرسش که این مصالحه تا چه حد به هندسه‌ی میله‌ای وابسته است، تنها با مقایسه با یک هندسه‌ی مرجع، یعنی نانوکره‌ی هم‌حجم، ممکن می‌شود.

\section{اهمیت و ضرورت تحقیق}
اهمیت کاربردی این تحقیق به طور مستقیم در توسعه، طراحی و بهینه‌سازی چشمه‌های تک‌فوتونی تحریک‌پالسی (\lr{On-demand Single-photon Sources}) نهفته است \cite{novotny,pelton}. این چشمه‌ها، قلب تپنده‌ی فناوری‌های نوین اطلاعات کوانتومی، از جمله مخابرات کوانتومی، محاسبات نوری خطی و شبکه‌های توزیع کلید کوانتومی (\lr{QKD}) مانند پروتکل \lr{BB84} به شمار می‌روند \cite{novotny}. در این سامانه‌های پیشرفته، امنیت تبادل داده و سرعت پردازش اطلاعات کاملاً وابسته به توانایی تولید دقیقاً «یک فوتون» در هر پالس تحریکی است. نرخ فروافت کل گسیلنده (که معکوس طول عمر حالت برانگیخته است)، مستقیماً بیشینه‌ی نرخ تکرار چشمه (سرعت ارسال داده) را تعیین می‌کند. کاهش این طول عمر از مقیاس نانوثانیه به پیکوثانیه به کمک اثر پورسل، گامی مهم در افزایش پهنای باند مخابرات کوانتومی است \cite{novotny,maier}.

از سوی دیگر، برای تولید فوتون‌های کاملاً تمایزناپذیر (\lr{Indistinguishable photons}) که شرط اساسی برای انجام عملیات‌های تداخل‌سنجی کوانتومی (مانند اثر هنگ-او-مندل) است، طول عمر تابشی گسیلنده باید در مقایسه با زمان وافازی (\lr{Dephasing time}) در محیط جامد بسیار کوتاه باشد \cite{novotny,pelton}. با تسریع نرخ گسیل از طریق نانوآنتن پلاسمونی، فوتون ساطع‌شده پیش از آنکه تحت تأثیر نوسانات فونونی محیط دچار وافازی (\lr{Dephasing}) شود، از سامانه خارج می‌گردد \cite{maier,pelton}. سنجش و تحلیل دقیق سهم کانال‌های تابشی در این پژوهش مشخص می‌کند که یک چشمه (مثلاً یک نقطه‌ی کوانتومی کلوئیدی یا یک مرکز تهی‌جای نیتروژن، \lr{NV}، در الماس) در چه فاصله‌ای از نانوآنتن طلا می‌تواند بیشترین فوتون قابل جمع‌آوری (بالاترین روشنایی یا \lr{Brightness}) را تولید کند، بدون آنکه در دام اتلاف اهمی فلز گرفتار شود \cite{novotny,anger}.

علاوه بر جنبه‌های کاربردی گسترده، این پژوهش از دیدگاه نظری و بنیادین نیز در فیزیک نانواپتیک اهمیت دارد \cite{novotny}. زبانی که امروزه در ادبیات علمی برای توصیف گسیل خودبه‌خودی و مدهای الکترومغناطیسی به کار می‌رود، عمدتاً از فیزیک اتمی به ارث رسیده است. در این زبان کلاسیک، برچسب‌ها و اعداد کوانتومی (مانند اعداد سمتی $l$ و $m$) تماماً بر پایه‌ی «تقارن کروی» بنا شده‌اند \cite{novotny}. نانوکره‌ی طلا دقیقاً در همین چارچوب قرار می‌گیرد و مسئله‌ی گسیلنده در کنار آن، به کمک نظریه‌ی می (\lr{Mie}) و بسط چندقطبی، به صورت تحلیلی و دقیق قابل حل است \cite{mie,ruppin,gersten}. با گذر از هندسه‌ی کروی متقارن به سمت یک نانومیله، تقارن کروی پیوسته‌ی ($SO(3)$) به شدت شکسته شده و سامانه به یک هندسه با تقارن استوانه‌ای ($C_{\infty v}$) تنزل می‌یابد \cite{maier}. این افت تقارن، رفتار الکترودینامیکی سامانه را کاملاً دگرگون ساخته و باعث می‌شود نانوذره دارای دو تشدید مجزای طولی و عرضی با شدت، الگوی تابشی و انرژی کاملاً متفاوت باشد \cite{maier}. از این رو، ضروری است که مدل‌های همسانگرد متعارف به چالش کشیده شوند و دامنه‌ی اعتبار رویکردهای تک‌پارامتری در ساختارهای شدیداً ناهمسانگردی چون نانومیله‌های طلا، به طور دقیق و کمّی، با استفاده از ابزارهای حل عددی و در مقایسه با حل تحلیلی نانوکره، ارزیابی گردند \cite{novotny,koenderink}.

\section{پیشینه‌ی پژوهش}
تاریخچه‌ی کنترل و دستکاری گسیل خودبه‌خودی با مقاله‌ی بنیادین و تاریخی ادوارد میلز پورسل در سال ۱۹۴۶ آغاز شد \cite{purcell}. پورسل برای نخستین بار پیش‌بینی کرد که احتمال گذار یک سامانه (در بحث او، گذارهای اسپین هسته‌ای در بسامدهای رادیویی) با قرارگیری در یک مدار تشدیدی تنظیم‌شده، نسبت به فضای آزاد به شدت تغییر می‌کند \cite{purcell}. تا دهه‌ی ۱۹۷۰، این پدیده عمدتاً یک مفهوم نظری در فیزیک امواج مایکروویو باقی ماند، تا آنکه کارل هاینتس درکسهاگه (\lr{K. H. Drexhage}) با انجام آزمایش‌های کلاسیک و هوشمندانه بر روی تک‌لایه‌هایی از مولکول‌های فلورسنت آلی در فواصل کنترل‌شده از یک آینه‌ی فلزی، توانست تغییرات نوسانی در طول عمر فلورسانس را به صورت تجربی مشاهده کند \cite{drexhage,novotny}. این نوسانات، تأییدی تجربی بر تداخل میدان گسیل‌شده با میدان بازتاب‌یافته از سطح فلز و تغییر چگالی حالت‌های موضعی (\lr{LDOS}) بود \cite{novotny,drexhage}.

در سال ۱۹۸۷، الی یابلونوویچ (\lr{Eli Yablonovitch}) و ساجیو جان (\lr{Sajeev John}) به صورت مستقل ایده‌ی «بلورهای فوتونی» (\lr{Photonic Crystals}) را مطرح کردند \cite{yablonovitch,john}؛ ساختارهای متناوب دی‌الکتریکی که می‌توانستند با ایجاد یک «گاف باند نوری» (\lr{Photonic Bandgap})، چگالی حالت‌های فوتونی را در یک محدوده‌ی بسامدی خاص به صفر رسانده و در نتیجه، گسیل خودبه‌خودی را کاملاً مهار کنند \cite{yablonovitch,john,novotny}. با وجود دستاوردهای مهم این حوزه در دهه‌های ۸۰ و ۹۰ میلادی، حجم مدی در این ساختارهای تمام‌دی‌الکتریک هیچ‌گاه نتوانست به مقادیر زیر طول موج برسد، چرا که قوانین پراش همواره بر این سامانه‌ها حاکم بود \cite{novotny,maier}.

هم‌زمان، نظریه‌ی فروافت گسیلنده در کنار یک کره‌ی فلزی، که ساده‌ترین نانوآنتن ممکن است، در اوایل دهه‌ی ۱۹۸۰ توسط گرستن و نیتزان (\lr{Gersten and Nitzan}) و نیز روپین (\lr{Ruppin}) به صورت تحلیلی تکمیل شد \cite{gersten,ruppin}. این کارها نشان دادند که در کنار یک کره‌ی فلزی کوچک، نرخ کل فروافت با کاهش فاصله به شدت افزایش می‌یابد، اما سهم غیرتابشی آن به صورت $1/d^3$ رشد کرده و در فواصل کوتاه بر تابش غالب می‌شود؛ نتیجه‌ای که بعدها با محاسبات دقیق نظریه‌ی می نیز تأیید شد \cite{mertens}.

انقلاب اصلی در تمرکز میدان الکترومغناطیسی با ظهور علم نانوفوتونیک و پلاسمونیک در اواخر دهه‌ی ۱۹۹۰ و اوایل ۲۰۰۰ رخ داد \cite{maier,novotny}. ویلیام بارنز (\lr{W. L. Barnes}) و توماس ابسن (\lr{T. W. Ebbesen}) در سال ۲۰۰۳ در مقاله‌ای مروری در مجله‌ی \lr{Nature} نشان دادند که پلاسمون‌های سطحی می‌توانند نور را در ابعاد بسیار کوچک‌تر از طول موج در فصل مشترک فلز و دی‌الکتریک محبوس کنند و مسیری جدید برای کوچک‌سازی مدارهای نوری ارائه دهند \cite{barnes,maier}. در نیمه‌ی دوم دهه‌ی ۲۰۰۰، لوکاس نووتنی (\lr{Lukas Novotny}) و همکارانش مفهوم «نانوآنتن‌های نوری» را با الهام از آنتن‌های رادیویی، به دنیای فیزیک نور معرفی کردند \cite{novotny,bharadwaj}.

در سال ۲۰۰۶، دو آزمایش پیشگامانه توسط گروه آنگر (\lr{Anger}) و بهارادواج (\lr{Bharadwaj}) و نیز گروه کوهن (\lr{Kühn}) و صندوق‌دار انجام پذیرفت \cite{anger,kuhn,novotny}. آن‌ها با استفاده از میکروسکوپ پروب روبشی (\lr{SPM}) و اتصال یک نانوذره‌ی «کروی» طلا (به قطر حدود ۸۰ تا ۱۰۰ نانومتر) به نوک پروب، یک تک‌مولکول فلورسنت را در فواصل مختلف از سطح نانوذره روبش کردند \cite{anger,kuhn}. نتایج این آزمایش‌ها پرده از یک رفتار دوگانه برداشت: با نزدیک کردن مولکول به نانوذره، ابتدا شدت فلورسانس به دلیل تقویت میدان نزدیک به شدت افزایش می‌یافت، اما با کاهش فاصله به کمتر از ۵ نانومتر، پدیده‌ی خاموشی غیرتابشی (\lr{Quenching}) بر سامانه چیره شده و تابش مفید سامانه به صفر میل می‌کرد \cite{anger,novotny}. بیشینه‌ی تقویت فلورسانس در این آزمایش‌ها حداکثر در حد یک مرتبه‌ی بزرگی (چند تا چند ده برابر) بود \cite{anger,kuhn}. این مشاهدات، مرز فیزیکی کارایی نانوآنتن‌ها را در سطح تک‌مولکول به روشنی نشان داد \cite{novotny,maier}.

در سال ۲۰۰۷، لاوینیا روگوبته (\lr{Rogobete}) و وحید صندوق‌دار (\lr{Vahid Sandoghdar}) با استفاده از شبیه‌سازی‌های عددی دقیق نشان دادند که نانوآنتن‌های فلزی پتانسیل آن را دارند که نرخ گسیل تک‌مولکول‌ها را صدها تا هزاران برابر افزایش دهند، به شرط آنکه گسیلنده در نقاط داغ ساختارهای تشدیدی قرار گیرد \cite{rogobete,novotny}. در همان سال، مرتنز (\lr{Mertens})، کوئندرینک و پولمن نشان دادند که مدل ساده‌شده‌ی گرستن-نیتزان، پس از اصلاح، با نظریه‌ی دقیق برای نانوکره‌های فلزی سازگاری خوبی دارد \cite{mertens}. در سال ۲۰۱۰، آلبرت کوئندرینک (\lr{A. F. Koenderink}) در مقاله‌ای انتقادی و ساختارگشا در مجله \lr{Optics Letters} بحث کرد که استفاده از پارامتر «فاکتور پورسل» به تنهایی و بدون تفکیک کانال‌های فروافت برای ارزیابی نانوآنتن‌های پلاسمونی، به دلیل ماهیت شدیداً اتلافی فلزات، گمراه‌کننده است \cite{koenderink}. وی نشان داد که افزایش نرخ کل فروافت لزوماً به معنای تولید فوتون‌های بیشتر نیست و محققان باید «بازده کوانتومی آنتن» و نرخ تابشی را به صورت مجزا محاسبه و گزارش نمایند \cite{koenderink,novotny}.

در دهه‌ی گذشته و به خصوص از سال ۲۰۱۱ به بعد، با پیشرفت روش‌های لیتوگرافی پرتو الکترونی (\lr{EBL}) و سنتزهای شیمیایی، تمرکز جامعه‌ی علمی از نانوکره‌های متقارن به سمت نانوآنتن‌های با تقارن شکسته تغییر یافت \cite{maier,giannini}. یکی از دلایل این تغییر، پیش‌تر توسط زونیکسن (\lr{Sönnichsen}) و همکاران در سال ۲۰۰۲ نشان داده شده بود: میرایی پلاسمون در نانومیله‌های طلا، به دلیل دور شدن تشدید طولی از لبه‌ی گذارهای بین‌نواری، به مراتب کمتر از نانوکره‌های طلاست \cite{sonnichsen}. محققانی مانند وینچنزو جیانینی (\lr{Giannini}) و استفان مایر (\lr{Stefan Maier}) نشان دادند که نانوآنتن‌های پاپیونی (\lr{Bowtie}) و نانومیله‌ها، بستر بسیار مناسب‌تری برای کنترل گسیل فراهم می‌آورند \cite{giannini,maier}. در سال ۲۰۱۴، خاطوا (\lr{Khatua}) و همکاران به صورت تجربی نشان دادند که یک تک‌نانومیله‌ی طلا می‌تواند فلورسانس یک تک‌مولکول را تا حدود هزار برابر تقویت کند \cite{khatua}؛ مقداری که با تقویت حداکثر چند ده برابری نانوکره‌ها در آزمایش‌های سال ۲۰۰۶ تفاوت چشمگیری دارد. در سال ۲۰۱۵، متیو پلتون (\lr{Matthew Pelton}) در مروری جامع در مجله‌ی \lr{Nature Photonics} جمع‌بندی کرد که نانومیله‌های طلا به دلیل دارا بودن تشدید طولی قدرتمند و ایجاد اثر برق‌گیر در نوک‌های خود، می‌توانند جفت‌شدگی بسیار کارآمدتری نسبت به کره‌های هم‌حجم ایجاد کنند \cite{pelton,novotny}.

امروزه حل‌کننده‌های عددی تمام‌موج مانند \lr{FDTD} ابزار استاندارد طراحی نانوآنتن‌ها هستند \cite{taflove}. با این وجود، رفتار دقیق و کمّی نانومیله‌های کلاهک‌دار (\lr{Hemispherically-capped}) در فواصل بحرانی زیر ۱۰ نانومتر و چگونگی تفکیک دقیق کانال‌های تابشی از اتلاف اهمی طلا، و به ویژه مقایسه‌ی کمّی آن با نانوکره‌ی هم‌حجم در شرایط یکسان، همچنان نیازمند بررسی‌های سه‌بعدی با دقت بالاست \cite{novotny,maier}. پژوهش حاضر، با آگاهی از این سیر تکاملی، بر روی پر کردن این خلأ محاسباتی در رژیم میدان نزدیک و یافتن نقطه‌ی بهینه‌ی کارکرد چشمه‌های تک‌فوتونیِ جفت‌شده با نانومیله‌ی طلا متمرکز است.

\section{اهداف پژوهش}
بر مبنای ضرورت‌های علمی و چالش‌های بنیادین مطرح‌شده در بیان مسئله و خلأهای شناسایی‌شده در پیشینه‌ی پژوهش، اهداف کلان و خرد این پایان‌نامه به شرح زیر تدوین شده است:

\begin{enumerate}
\item \textbf{توسعه و پیاده‌سازی مدل الکترومغناطیسی با دقت بالا:} هدف نخست، ایجاد یک بستر شبیه‌سازی سه‌بعدی دقیق مبتنی بر روش تفاضل محدود در حوزه‌ی زمان (\lr{3D-FDTD}) است \cite{taflove}. این مدل باید قادر باشد با استفاده از لایه‌های جاذب کاملاً تطبیق‌یافته (\lr{PML}) \cite{berenger} و مش‌بندی موضعی بسیار ریز (در حد یک نانومتر)، رفتار الکترودینامیکی یک نانومیله‌ی طلا را با استفاده از مدل پاشندگی وابسته به بسامد \cite{maier} و داده‌های تجربی جانسون-کریستی \cite{jc} به درستی ردگیری نماید.

\item \textbf{تفکیک نرخ‌های گذار و محاسبه‌ی بازده کوانتومی:} برخلاف رویکردهایی که تنها نرخ کل فروافت (فاکتور پورسل) را گزارش می‌کنند، این پژوهش قصد دارد با تعریف استراتژیک جعبه‌های مانیتورینگ انتقال توان (\lr{Transmission boxes}) در نرم‌افزار، کل انرژی خارج‌شده از دوقطبی کلاسیک را به دو بخش «توان تابشی به فضای دوردست» و «توان اتلافی در شبکه‌ی فلز» تفکیک نماید تا بازده کوانتومی ظاهری سامانه به طور قطعی مشخص گردد \cite{novotny,koenderink}.

\item \textbf{پویش سیستماتیک فواصل (\lr{Gap Sweep}) برای یافتن نقطه‌ی بهینه:} بررسی دقیق رفتار سامانه هنگام تغییر فاصله‌ی گسیلنده از سطح فلز (از ۳ تا ۲۰ نانومتر). هدف این بخش، رصد کردن پدیده‌ی خاموشی غیرتابشی در فواصل کوتاه و یافتن مصالحه‌ی فیزیکی میان بالاترین فاکتور پورسل ممکن و بیشترین بازده تابشی است تا «گاف طلایی» برای قرارگیری چشمه‌های فوتونی به دست آید \cite{novotny,anger}. نتایج این پویش با رفتار یک نانوکره‌ی طلای هم‌حجم در فواصل یکسان، که حل تحلیلی دقیق دارد، مقایسه می‌شود \cite{ruppin,gersten}.

\item \textbf{کمّی‌سازی میزان ناهمسانگردی جهتی (\lr{Anisotropy}):} هدف نهایی، مقایسه‌ی سیستماتیک پاسخ ساختار به جهت‌گیری‌های مختلف دوقطبی (موازی و عمود بر محور طولی نانومیله) و مقایسه‌ی آن با نانوکره‌ی هم‌حجم (گروه کنترل) است تا میزان قطبش‌پذیری سامانه و قابلیت نانوآنتن طلا به عنوان یک «فیلتر قطبشی فوتون» به صورت عددی نشان داده شود \cite{novotny,maier}.
\end{enumerate}

% ======================================== فصل ۲ ======================================================
\chapter{مبانی پلاسمونیک و خواص نوری فلزات}

\section{فیزیک الکترودینامیک فلزات و مدل پلاسما (مدل درود)}
برهم‌کنش فلزات با میدان‌های الکترومغناطیسی در ابعاد نانو را می‌توان به طور دقیق و منسجم در چارچوب فیزیک کلاسیک و بر پایه‌ی معادلات ماکروسکوپیک ماکسول توصیف کرد \cite{maier,novotny}. در نانوساختارهای فلزی، چگالی حامل‌های بار آزاد به قدری بالاست (در مرتبه‌ی $10^{22}$ الکترون بر سانتی‌متر مکعب) که فاصله‌ی ترازهای انرژی الکترونی در مقایسه با انرژی حرارتی محیط ($k_BT$) بسیار ناچیز می‌شود؛ از این رو، بررسی الکترودینامیکی این سامانه‌ها با مدل‌های کلاسیک و نیمه‌کلاسیک کاملاً معتبر و پاسخگو است \cite{maier}.

در بسامدهای پایین (مانند امواج رادیویی و مایکروویو)، فلزات به عنوان رساناهای ایده‌آل (\lr{Perfect Electrical Conductors - PEC}) عمل کرده و با بازتاب کامل امواج، مانع از نفوذ میدان به درون خود می‌شوند \cite{novotny,maier}. در این حالت، میدان‌های الکتریکی استاتیک در فواصل بسیار کوتاهی (طول غربالگری توماس-فرمی) توسط الکترون‌های رسانش خنثی یا غربال می‌شوند. با این حال، در بسامدهای نوری (محدوده‌ی مرئی و فروسرخ)، الکترون‌ها به دلیل داشتن جرم و اینرسی ذاتی نمی‌توانند به صورت آنی به تغییرات فوق‌سریع میدان محرک (با بسامدهای مرتبه‌ی صدها تراهرتز) پاسخ دهند \cite{maier}. در نتیجه، غربالگری ناقص می‌شود و میدان الکترومغناطیسی تا عمقی که به آن «عمق پوسته» (\lr{Skin depth}) می‌گویند و برای طلا در ناحیه‌ی مرئی حدود ۲۰ تا ۳۰ نانومتر است، به درون فلز نفوذ می‌کند؛ عمقی که با ابعاد خود نانوذرات قابل مقایسه است \cite{maier,novotny}.

در این رژیم نوری، رفتار فلز با یک تابع دی‌الکتریک مختلط و وابسته به بسامد، یعنی $\epsilon(\omega)=\epsilon_1(\omega)+i\epsilon_2(\omega)$ توصیف می‌گردد که تمایز سنتی میان بارهای آزاد و مقید را در معادلات ماکسول کم‌رنگ می‌کند \cite{maier}. بخش حقیقی این تابع ($\epsilon_1$) میزان قطبش‌پذیری محیط و توانایی فلز در ذخیره‌ی انرژی الکتریکی را تعیین می‌کند، در حالی که بخش موهومی آن ($\epsilon_2$) که مستقیماً با رسانندگی جریان متناوب فلز در ارتباط است، بیانگر میزان جذب و اتلاف انرژی الکترومغناطیسی (اتلاف اهمی) در داخل فلز می‌باشد \cite{maier,novotny}.

از دیدگاه فیزیکی، خواص نوری فلزات در یک محدوده‌ی بسامدی وسیع (از بسامدهای پایین تا محدوده‌ی فروسرخ) با استفاده از \textbf{مدل پلاسما (مدل درود - \lr{Drude Model})} توضیح داده می‌شود \cite{maier,novotny}. در این مدل، فلز به صورت گازی از الکترون‌های آزاد در نظر گرفته می‌شود که در برابر یک پس‌زمینه‌ی ثابت از هسته‌های یونی مثبت حرکت می‌کنند \cite{maier}. معادله‌ی حرکت یک الکترون منفرد تحت تأثیر میدان الکتریکی هماهنگ $E(t)=E_0e^{-i\omega t}$ بر اساس قانون دوم نیوتن به صورت زیر نوشته می‌شود \cite{maier}:
\begin{equation}
m_e\frac{d^2x}{dt^2}+m_e\gamma\frac{dx}{dt}=-eE\,.
\end{equation}
در این معادله، $m_e$ جرم مؤثر الکترون، $e$ بار پایه‌ی الکترون و $\gamma$ بسامد برخورد (نرخ میرایی) است که نشان‌دهنده‌ی اصطکاک ناشی از برخورد الکترون‌ها با یون‌های شبکه می‌باشد \cite{maier}. با حل این معادله‌ی دیفرانسیل، جابه‌جایی الکترون ($x$) به دست می‌آید که از طریق آن می‌توان قطبش ماکروسکوپیک فلز ($P=-nex$) را محاسبه نمود ($n$ چگالی الکترون‌های آزاد است). با جایگذاری این قطبش در رابطه‌ی بردار جابه‌جایی الکتریکی ($D=\epsilon_0E+P$)، تابع دی‌الکتریک معروف درود به شکل زیر استخراج می‌گردد \cite{maier,novotny}:
\begin{equation}
\epsilon(\omega)=1-\frac{\omega_p^2}{\omega^2+i\gamma\omega}\,.
\end{equation}
در این رابطه‌ی بنیادین، $\omega_p$ نشان‌دهنده‌ی «بسامد پلاسما» است که به صورت $\omega_p=\sqrt{ne^2/(m_e\epsilon_0)}$ تعریف می‌شود \cite{maier}. بررسی این معادله نشان می‌دهد که اگر در بسامد خاصی (با صرف‌نظر از میرایی $\gamma$)، تابع دی‌الکتریک محیط صفر شود ($\epsilon(\omega)=0$)، نوسانات جمعی طولی در محیط رخ می‌دهد. کوانتای این نوسانات باری، «پلاسمون‌های حجمی» (\lr{Volume Plasmons}) نامیده می‌شوند \cite{maier}. بسامد پلاسما در واقع نشان‌دهنده‌ی بسامد طبیعی نوسانات دسته‌جمعی دریای الکترون‌های رسانش است. در بسامدهای کمتر از $\omega_p$ (که ناحیه‌ی مرئی و فروسرخ برای بیشتر فلزات نجیب را در بر می‌گیرد)، بخش حقیقی تابع دی‌الکتریک منفی است که خاصیت بازتاب‌دهندگی و پلاسمونی فلز را تضمین می‌کند \cite{maier,novotny}.

\section{گذارهای بین‌نواری، اتلاف اهمی و مدل درود-لورنتز}
ساختار داخلی و دینامیک الکترون‌ها در فلزات، تعیین‌کننده‌ی اصلی اتلاف و نحوه‌ی پاسخگویی پلاسما به تحریکات خارجی است \cite{maier}. در مدل درود، جمله‌ی میرایی ($\gamma=1/\tau$) ناشی از برخورد الکترون‌ها با شبکه (فونون‌ها)، ناخالصی‌ها و سایر الکترون‌هاست. این برخوردها منشأ اصلی اتلاف اهمی در فلزات به شمار می‌روند و دقیقاً همین اتلاف است که در فواصل چند نانومتری از سطح فلز، منجر به پدیده‌ی خاموشی غیرتابشی (\lr{Quenching}) و افت شدید بازده تابشی در گسیلنده‌های کوانتومی می‌شود \cite{novotny,anger}.

با این وجود، ساختار واقعی پلاسمون‌ها نیازمند اصلاحات کوانتومی است. با اعمال اصل طرد پائولی و آمار فرمی-دیراک در نظریه‌ی «درود-سامرفلد»، مشخص می‌شود که تنها الکترون‌های نزدیک به سطح فرمی در پاسخ نوری مشارکت کرده و با سرعت فرمی ($v_F$) حرکت می‌کنند \cite{maier,novotny}. این موضوع در نانوساختارها پیامدهای مهمی دارد؛ برای نانوذراتی که شعاع آن‌ها ($R$) از طول پویش آزاد الکترون‌ها ($l_{mfp}$ - \lr{Mean Free Path}) کوچک‌تر است، اثرات مرزی غالب می‌شود \cite{maier,kreibig}. در این حالت، برخورد الکترون‌ها با دیواره‌ها و سطح ذره به شدت افزایش یافته و نرخ میرایی را به صورت تابعی از ابعاد هندسی افزایش می‌دهد \cite{kreibig,maier}:
\begin{equation}
\gamma=\gamma_{\sub{bulk}}+A\frac{v_F}{R}\,,
\end{equation}
که در آن $\gamma_{\sub{bulk}}$ نرخ میرایی توده‌ی فلز و $A$ یک ثابت بدون‌بعد از مرتبه‌ی یک است. این پدیده نشان می‌دهد که اتلاف اهمی در ساختارهای نانومتری بسیار کوچک (زیر ۱۰ نانومتر) به مراتب شدیدتر از فلزات توده‌ای (\lr{Bulk}) است \cite{maier,kreibig}. برای نانومیله‌ی این پژوهش (شعاع مقطع ۱۰ نانومتر)، با $v_F\approx1.4\times10^{6}$ متر بر ثانیه، سهم سطحی $\hbar Av_F/R$ در حدود چند صدم الکترون‌ولت است؛ یعنی با میرایی توده‌ای طلا ($\hbar\gamma_{\sub{bulk}}\approx0.07$ الکترون‌ولت) قابل مقایسه است. برای نانوکره‌ی هم‌حجم (شعاع حدود ۱۶ نانومتر) این سهم کمتر است. از آنجا که در شبیه‌سازی‌ها از داده‌های توده‌ای جانسون-کریستی استفاده شده است، این اثر لحاظ نشده و در عمل قله‌های تشدید اندکی پهن‌تر و کوتاه‌تر، و بازده‌ها اندکی کمتر از مقادیر گزارش‌شده خواهند بود؛ این نکته به عنوان یکی از محدودیت‌های مدل در نظر گرفته می‌شود \cite{kreibig,jc}.

از سوی دیگر، ساختار پلاسمون در فلزات نجیب مانند طلا و نقره، پیچیدگی‌های بیشتری نسبت به مدل گاز الکترونی آزاد دارد \cite{maier}. پس‌زمینه‌ی مثبت یون‌ها در این فلزات کاملاً صلب نیست، بلکه دارای قطبش پسماند است که با ثابت دی‌الکتریک بسامد بالا ($\epsilon_\infty$) به مدل پایه اضافه می‌شود \cite{maier}. افزون بر این، در فلز طلا، نوار کاملاً پرشده‌ی $5d$ در فاصله‌ی بسیار نزدیکی (تنها حدود $2.4\,\mathrm{eV}$، معادل طول موج تابشی حدود ۵۱۷ نانومتر) از سطح فرمی قرار دارد \cite{maier,novotny,jc}. هنگامی که طلا تحت تابش فوتون‌هایی با انرژی بیشتر از این آستانه (در ناحیه‌ی سبز، آبی و فرابنفش) قرار می‌گیرد، الکترون‌ها انرژی فوتون را جذب کرده و از نوار $d$ به حالت‌های خالی نوار هدایت ($sp$) برانگیخته می‌شوند \cite{maier}. این «گذارهای بین‌نواری» (\lr{Interband Transitions}) منجر به ظهور یک مسیر جذب بسیار قدرتمند و اتلاف شدید انرژی در طول موج‌های مرئی کوتاه می‌گردند \cite{maier,novotny}. همان‌گونه که در ادامه خواهیم دید، تشدید نانوکره‌ی طلا دقیقاً در همین ناحیه‌ی پراتلاف قرار می‌گیرد، در حالی که تشدید طولی نانومیله از آن فاصله دارد \cite{link,sonnichsen}. مدل درود به تنهایی قادر به توصیف این پدیده‌ی کوانتومی نیست؛ از این رو، نوسانگرهای مقید لورنتز به معادله‌ی درود افزوده می‌شوند تا «مدل درود-لورنتز» شکل گیرد \cite{maier,novotny}:
\begin{equation}
\epsilon(\omega)=\epsilon_\infty-\frac{\omega_p^2}{\omega(\omega+i\gamma)}+\sum_{j=1}^{M}\frac{f_j\,\omega_{0,j}^2}{\omega_{0,j}^2-\omega^2-i\Gamma_j\omega}\,.
\end{equation}
در این رابطه، $f_j$ قدرت نوسانگر، $\omega_{0,j}$ بسامد تشدید گذار بین‌نواری و $\Gamma_j$ پهنای باند میرایی آن است \cite{maier,novotny}. استفاده از این مدل بسط‌یافته و تطبیق آن با داده‌های تجربی دقیق (مانند داده‌های جانسون و کریستی \cite{jc})، برای مدل‌سازی دقیق ساختارهای پلاسمونی طلا و درک صحیح خاموشی غیرتابشی الزامی است \cite{maier,novotny}.

در نهایت، با وجود شناخت ساختار پلاسمون‌های حجمی، به دلیل ماهیت کاملاً طولی این برانگیختگی‌ها (یعنی بردار میدان الکتریکی موازی با بردار انتشار موج است)، آن‌ها بر اساس قوانین پایه‌ی الکترودینامیک به هیچ وجه نمی‌توانند با امواج الکترومغناطیسی عرضی نور فضای آزاد جفت شوند \cite{maier}. این محدودیت فیزیکی، ضرورت استفاده از مرزهای هندسی و «پلاسمون‌های سطحی» را برای محبوس‌سازی و هدایت امواج نوری ایجاب می‌کند \cite{maier,barnes}.

\section{پلاسمون پلاریتون سطحی (\lr{SPP}) و فیزیک میدان نزدیک}
همان‌طور که پیش‌تر اشاره شد، پلاسمون‌های حجمی به دلیل ماهیت طولی خود قادر به جفت‌شدگی با امواج الکترومغناطیسی عرضی (فوتون‌ها) نیستند \cite{maier}. این مانع بنیادین با انتقال به مرز مشترک یک فلز (با تابع دی‌الکتریک $\epsilon_m$) و یک عایق دی‌الکتریک (با تابع دی‌الکتریک $\epsilon_d$) برطرف می‌شود \cite{maier}. «پلاسمون پلاریتون‌های سطحی» (\lr{Surface Plasmon Polaritons - SPP}) برانگیختگی‌های الکترومغناطیسی مقیدی هستند که در امتداد این مرز مشترک منتشر شده و نتیجه‌ی جفت‌شدگی مستقیم نور با نوسانات باری سطح فلز می‌باشند \cite{maier,novotny,barnes}.

برای استخراج روابط حاکم بر \lr{SPP}، حل معادلات موج استخراج‌شده از معادلات ماکسول در هر دو محیط ضروری است \cite{maier}. با فرض اینکه سطح مشترک در صفحه‌ی $z=0$ قرار دارد و موج در راستای $x$ منتشر می‌شود، میدان‌ها باید در راستای $z$ (عمود بر سطح) به صورت نمایی افت کنند (میدان ناپایا یا \lr{Evanescent}) \cite{maier,novotny}. با اعمال شرایط مرزی پیوستگی مؤلفه‌های مماس میدان الکتریکی ($E_x$) و مؤلفه‌های عمودی بردار جابه‌جایی ($D_z$)، ثابت می‌شود که این امواج مقید تنها با \textbf{قطبش مغناطیسی عرضی (\lr{TM}) یا (\lr{p-polarized})} می‌توانند وجود داشته باشند \cite{maier,novotny}. به عبارت دیگر، وجود مؤلفه‌ی میدان الکتریکی عمود بر سطح فلز برای تحریک این بارهای سطحی و ایجاد مدهای پلاسمونی کاملاً الزامی است؛ امواج \lr{TE} هیچ‌گونه پاسخ سطحی مقیدی تولید نمی‌کنند \cite{maier}.

با جایگذاری پاسخ‌ها در معادله‌ی موج و اعمال شرایط مرزی، بردار موج این برانگیختگی‌ها در امتداد سطح انتشار ($k_{spp}$) به صورت رابطه‌ی پاشندگی زیر به دست می‌آید \cite{maier,novotny}:
\begin{equation}
k_{spp}=k_0\sqrt{\frac{\epsilon_m\epsilon_d}{\epsilon_m+\epsilon_d}}\,,
\end{equation}
که در آن $k_0=\omega/c$ بردار موج نور در فضای آزاد است \cite{maier}. شرط فیزیکی و ریاضی برای واقعی بودن این بردار موج (و در نتیجه امکان انتشار موج روی سطح) این است که بخش حقیقی تابع دی‌الکتریک فلز نه تنها باید منفی باشد، بلکه اندازه‌ی (قدر مطلق) آن باید از ثابت دی‌الکتریک محیط مجاور فراتر رود (یعنی $\epsilon_1<-\epsilon_d$) \cite{maier,novotny}. فلزات نجیب مانند طلا، این شرط حیاتی را در رژیم‌های طول موج مرئی و فروسرخ نزدیک به خوبی برآورده می‌کنند و میزبان مناسبی برای امواج \lr{SPP} به شمار می‌روند \cite{maier}.

تحلیل منحنی پاشندگی \lr{SPP} نشان می‌دهد که تکانه‌ی این امواج سطحی ($\hbar k_{spp}$) همواره از تکانه‌ی فوتون‌های فضای آزاد ($\hbar k_0$) در همان بسامد بزرگ‌تر است \cite{maier,novotny}. منحنی پاشندگی \lr{SPP} همواره در زیر «خط نور» (\lr{Light line}) قرار می‌گیرد. این «عدم تطابق تکانه» (\lr{Momentum mismatch}) سبب می‌شود که پلاسمون‌های سطحی بر روی یک سطح تخت فلزی به هیچ‌وجه نتوانند به طور مستقیم توسط نور تابشی برانگیخته شوند (زیرا نور فرودی تکانه‌ی کافی برای ایجاد \lr{SPP} را ندارد) \cite{maier,novotny}. غلبه بر این محدودیت فیزیکی، نیازمند استفاده از منشورها (پیکربندی کرشمان یا اوتو)، توری‌های پراش، و یا ساختارهای هندسی خاص با ابعاد محدود مانند «نانوذرات» است تا با شکستن تقارن انتقالی در مرزها، تکانه‌ی اضافی مورد نیاز برای جفت‌شدگی نور با پلاسما تأمین گردد \cite{maier,novotny}.

میدان الکترومغناطیسی \lr{SPP} به سطح فلز چسبیده است و دامنه‌ی آن در راستای $z$ افت نمایی دارد \cite{maier}. عمق نفوذ یا «طول افت» (\lr{Decay length}) آن در داخل دی‌الکتریک، حجم مؤثر ناحیه‌ی نزدیک (\lr{Near-field volume}) را با دقت تعیین می‌کند \cite{novotny,maier}. این پارامتر در الکترودینامیک کوانتومی بسیار حائز اهمیت است؛ یک گسیلنده‌ی کوانتومی (مانند یک دوقطبی نقطه‌ای) تنها در صورتی می‌تواند انرژی خود را با بالاترین راندمان به مدهای پلاسمونی منتقل کند که در درون همین میدان موضعی و میرا قرار گرفته باشد \cite{novotny,maier}.

با این وجود، دینامیک پلاسمون‌ها در فلزات واقعی همواره با چالش اتلاف همراه است \cite{maier}. وجود بخش موهومی در تابع دی‌الکتریک طلا ($\epsilon_2$) باعث می‌شود که بردار موج \lr{SPP} مختلط گردد ($k_{spp}=k'_{spp}+ik''_{spp}$) \cite{maier}. بخش موهومی این بردار ($k''_{spp}$)، افت نمایی انرژی موج را در طول حرکت روی سطح تعیین کرده و «طول انتشار» پلاسمون ($L_{spp}=1/2k''_{spp}$) را محدود می‌سازد \cite{maier,novotny}. در فیزیک نانواپتیک، یک مصالحه‌ی (\lr{Trade-off}) ذاتی میان «محبوس‌سازی میدان» و «اتلاف انرژی» وجود دارد \cite{maier,novotny}. هرچه میدان الکترومغناطیسی شدیدتر در مجاورت سطح محبوس شود (افزایش $k'_{spp}$)، نفوذ میدان به درون فلز به ناچار افزایش می‌یابد \cite{maier}. این نفوذ بیشتر به معنای درگیر شدن میدان با گذارهای بین‌نواری و پراکندگی‌های فونونی قدرتمند طلا است که اتلاف اهمی را به شدت بالا می‌برد \cite{maier,novotny}. همین سازوکار فیزیکی است که موجب می‌شود با کاهش بیش از حد فاصله‌ی گسیلنده از سطح فلز، فاکتور پورسل به شدت بالا رود، اما انرژی آزادشده به جای تبدیل شدن به فوتون‌های دوردست، مستقیماً صرف تغذیه‌ی مدهای پلاسمونی میرا و گرم کردن ساختار گردد \cite{novotny,anger}.

\section{تشدید پلاسمون سطحی جایگزیده (\lr{LSPR}) و مدل‌های هندسی}
در شرایطی که سطح مشترک دی‌الکتریک و فلز دیگر یک صفحه‌ی تخت نامتناهی نیست و فلز به صورت یک نانوذره با ابعادی بسیار کوچک‌تر از طول موج نور تابشی شکل داده شده است ($2a\ll\lambda$)، فاز میدان الکترومغناطیسی فرودی در سراسر حجم ذره تقریباً ثابت فرض می‌شود \cite{maier,novotny}. در این رژیم که در الکترودینامیک به «تقریب شبه‌استاتیک» (\lr{Quasi-static approximation}) معروف است، مسئله‌ی پیچیده‌ی انتشار امواج به یک مسئله‌ی الکترواستاتیک ساده‌تر تقلیل می‌یابد \cite{maier,novotny}.

بر اساس این تقریب، برای یک نانوکره‌ی فلزی به شعاع $a$ که در یک میدان الکتریکی یکنواخت استاتیک قرار گرفته است، می‌توان پتانسیل الکتریکی ($\Phi$) را از طریق حل معادله‌ی لاپلاس ($\nabla^2\Phi=0$) محاسبه نمود \cite{maier,novotny}. با اعمال شرایط مرزی (پیوستگی پتانسیل و مؤلفه‌ی عمودی جابه‌جایی الکتریکی) در سطح کره، پتانسیل در فضای خارج ذره به دست می‌آید که معادل پتانسیل یک دوقطبی ایده‌آل قرارگرفته در مرکز کره است. از این طریق، «قطبش‌پذیری دوقطبی» (\lr{Dipole Polarizability - $\alpha$}) برای نانوکره استخراج می‌گردد \cite{maier,novotny,bohren}:
\begin{equation}\label{eq:alpha-sphere}
\alpha=4\pi a^3\,\frac{\epsilon_m-\epsilon_d}{\epsilon_m+2\epsilon_d}\,,
\end{equation}
که در آن گشتاور دوقطبی القایی برابر $\mathbf{p}=\epsilon_0\epsilon_d\,\alpha\,\mathbf{E}_0$ است و $\alpha$ بُعد حجم دارد \cite{maier}. میدان الکتریکی فرودی به عنوان نیروی محرک، کل ابر الکترونی نانوذره را نسبت به شبکه‌ی ثابت یون‌های مثبت جابه‌جا می‌کند؛ بارهای سطحی حاصل، یک نیروی بازگرداننده‌ی کولنی ایجاد می‌کنند و نتیجه، یک نوسان دوقطبی هماهنگ است \cite{maier}. بر اساس رابطه‌ی قطبش‌پذیری، مشخص است که هرگاه مخرج کسر به سمت صفر میل کند، یعنی بخش حقیقی تابع دی‌الکتریک فلز شرط $\epsilon_1=-2\epsilon_d$ (معروف به \textbf{شرط فرولیش} - \lr{Fröhlich Condition}) را برآورده سازد، دامنه‌ی این قطبش‌پذیری و نوسانات به شدت افزایش می‌یابد \cite{maier,novotny}. در واقعیت، به دلیل وجود $\epsilon_2$، دامنه محدود اما بسیار بزرگ می‌گردد. در این حالت، میدان الکترومغناطیسی فرودی به شدت توسط ذره جذب و پراکنده شده و میدان‌های موضعی در مجاورت سطح ذره به شدت تقویت می‌شوند؛ پدیده‌ای که موسوم به «تشدید پلاسمون سطحی جایگزیده» (\lr{Localized Surface Plasmon Resonance - LSPR}) است \cite{maier,novotny}.

برای نانوکره‌ی طلا در خلأ یا هوا ($\epsilon_d=1$)، شرط فرولیش در حدود ۵۰۰ تا ۵۳۰ نانومتر برقرار می‌شود؛ ناحیه‌ای که بر لبه‌ی گذارهای بین‌نواری طلا (حدود ۵۱۷ نانومتر) منطبق است \cite{kreibig,link,jc}. به همین دلیل، در این طول‌موج $\epsilon_2$ بزرگ است و تشدید نانوکره‌ی طلا پهن و پراتلاف است. نکته‌ی مهم‌تر آنکه چون رابطه‌ی (\ref{eq:alpha-sphere}) به شکل ذره وابسته نیست و اندازه‌ی ذره تنها از طریق اثرات تأخیر و میرایی سطحی وارد می‌شود، طول‌موج تشدید نانوکره تنها به مقدار اندکی (چند ده نانومتر) قابل جابه‌جایی است \cite{kreibig,bohren}.

برخلاف پلاسمون‌های سطحی منتشرشونده (\lr{SPP}) که به مرزهای تخت محدودند، در مدهای جایگزیده (\lr{LSPR})، شرایط تشدید به شدت تحت تأثیر هندسه‌ی ذره قرار دارد؛ زیرا شکل مرزها، توزیع بارها و نیروی بازگرداننده وارد بر ابر الکترونی را مستقیماً دیکته می‌کند \cite{maier,novotny}. هنگامی که تقارن کروی شکسته شده و ذره به سمت ساختارهای کشیده مانند بیضی‌گون‌ها و نانومیله‌های طلا میل می‌کند، پاسخ پلاسمونی رفتار کاملاً متفاوتی از خود نشان می‌دهد \cite{maier,giannini}.

برای یک بیضی‌گون یا یک نانومیله، قطبش‌پذیری سامانه دیگر یک اسکالر ساده نیست، بلکه یک تانسور است \cite{maier,novotny}. عدم تقارن در طول محورهای سه‌گانه ($x,y,z$) سبب می‌شود شرط فرولیش برای هر محور توسط یک پارامتر هندسی مستقل به نام «فاکتور واقطبش» (\lr{Depolarization factor}، $L_i$) اصلاح گردد \cite{maier,novotny,bohren}:
\begin{equation}\label{eq:alpha-i}
\alpha_i=V\,\frac{\epsilon_m-\epsilon_d}{\epsilon_d+L_i(\epsilon_m-\epsilon_d)}\,,
\end{equation}
که در آن $V$ حجم ذره و $\sum_i L_i=1$ است \cite{maier,bohren}. باید تأکید کرد که این رابطه و فاکتورهای واقطبش $L_i$ به صورت دقیق تنها برای بیضی‌گون‌ها (و حالت خاص آن‌ها، کره‌وارها) تعریف می‌شوند؛ زیرا تنها در بیضی‌گون است که میدان داخل ذره در یک میدان خارجی یکنواخت، یکنواخت باقی می‌ماند \cite{bohren}. برای استوانه‌ی کلاهک‌دار مورد مطالعه در این پژوهش، میدان داخلی یکنواخت نیست و $L_i$ تعریف تحلیلی دقیقی ندارد؛ از این رو، نانومیله در این بخش با یک کره‌وار کشیده‌ی هم‌ارز (\lr{Prolate spheroid}) تقریب زده می‌شود. این تقریب تنها برای توضیح کیفی رفتار نانومیله (رفع تبهگنی مدها و انتقال به سرخ مد طولی) به کار می‌رود و مقادیر دقیق، از حل عددی (\lr{FDTD}) در فصل پنجم به دست می‌آیند. برای کره ($L_x=L_y=L_z=1/3$)، این رابطه دقیقاً به رابطه‌ی (\ref{eq:alpha-sphere}) تبدیل می‌شود؛ بنابراین نانوکره حالت خاص و همسانگرد همین مدل است. این اصلاح هندسی موجب رفع تبهگنی (\lr{Degeneracy}) مدهای نوسانی شده و تشدید دوقطبی نانومیله را به دو مُد کاملاً متمایز تجزیه می‌کند \cite{maier,novotny}:
\begin{enumerate}
\item \textbf{مد عرضی (\lr{Transverse mode}):} مربوط به نوسان الکترون‌ها در امتداد محور کوتاه نانومیله. بسامد این مد معمولاً نزدیک به بسامد تشدید کره‌ی طلا (حدود ۵۲۰ نانومتر) است و به دلیل هم‌پوشانی با گذارهای بین‌نواری، به شدت اتلافی است \cite{maier,novotny,link}.
\item \textbf{مد طولی (\lr{Longitudinal mode}):} مربوط به نوسان بارهای الکترونی در امتداد محور اصلی و بلند نانومیله. با افزایش نسبت طول به عرض (\lr{Aspect Ratio})، فاکتور واقطبش طولی ($L_1$) کاهش یافته و بسامد این مد به شدت دچار «انتقال به سرخ» (\lr{Red-shift}) می‌شود \cite{maier,novotny,link}.
\end{enumerate}
این ویژگی به مهندسان نانوفوتونیک اجازه می‌دهد تا با تنظیم دقیق هندسه‌ی نانومیله، بسامد مد طولی را به دقت در محدوده‌ی طول موج‌های نوری تنظیم کنند تا با بسامد تابش یک گسیلنده‌ی کوانتومی خاص کاملاً جفت (\lr{Couple}) شود \cite{novotny,maier}. افزون بر این، چون تشدید طولی از لبه‌ی گذارهای بین‌نواری فاصله می‌گیرد، میرایی آن به مراتب کمتر از تشدید نانوکره است؛ زونیکسن و همکاران نشان دادند که پهنای خط تشدید نانومیله‌های طلا به طور چشمگیری باریک‌تر از نانوکره‌های طلاست \cite{sonnichsen}. علاوه بر این، در نانومیله‌ها، بارهای الکتریکی نوسانگر عمدتاً در نوک‌های ساختار تجمع می‌یابند. به دلیل انحنای شدید در قطب‌ها، خطوط میدان الکتریکی در این نواحی دچار تراکم فوق‌العاده‌ای شده و تقویت عظیم میدان نزدیک ایجاد می‌کنند (همان اثر برق‌گیر) \cite{novotny,bharadwaj}؛ ویژگی‌ای که در نانوکره به دلیل انحنای یکنواخت وجود ندارد.

با وجود این تقویت چشمگیر میدان، برهم‌کنش نانوذره با نور همواره با یک رقابت فیزیکی ذاتی میان دو سطح مقطع همراه است: «جذب» (\lr{Absorption}) و «پراکندگی» (\lr{Scattering}) \cite{maier,novotny}. بر اساس نظریه‌ی مای (\lr{Mie Theory}) \cite{mie,bohren} و بسط‌های شبه‌استاتیک، سطح مقطع پراکندگی متناسب با توان دوم قطبش‌پذیری (و در نتیجه متناسب با مربع حجم، $V^2$) است، در حالی که سطح مقطع جذب مستقیماً با بخش موهومی قطبش‌پذیری (متناسب با حجم، $V$) ارتباط دارد \cite{maier,novotny}:
\begin{equation}
C_{\sub{abs}}=k\,\Imag(\alpha)\quad,\quad C_{\sub{sca}}=\frac{k^4}{6\pi}|\alpha|^2\,,
\end{equation}
که در آن $k$ عدد موج در محیط اطراف است. نسبت این دو سطح مقطع با $V/\lambda^3$ متناسب است؛ بنابراین برای نانوذرات بسیار کوچک‌تر از طول موج، جذب اهمی بر پراکندگی غالب است و ذره انرژی را بیشتر جذب می‌کند تا پراکنده کند \cite{maier,novotny}. اما با بهینه‌سازی و افزایش ابعاد ذره تا یک حد مشخص، سهم پراکندگی تابشی بر جذب پیشی گرفته و نانومیله می‌تواند نقش یک آنتن کارآمد را ایفا کند \cite{novotny,maier}. باید توجه داشت که با بزرگ‌تر شدن ابعاد ذره، تقریب شبه‌استاتیک به دلیل بروز اثرات فیزیکی مانند «تأخیر فاز» (\lr{Phase Retardation}) و «میرایی تابشی» (\lr{Radiation Damping}) کارایی خود را از دست می‌دهد \cite{maier,novotny}. برای نانوکره، این اثرات با نظریه‌ی دقیق می قابل محاسبه‌اند \cite{mie,bohren}؛ اما در رژیم ابعادی واقعی و هندسه‌های کلاهک‌دار، حل تحلیلی دیگر پاسخگو نیست و استفاده از روش‌های حل عددی دقیق معادلات ماکسول، مانند روش تفاضل محدود در حوزه‌ی زمان (\lr{FDTD}) برای محاسبه‌ی دقیق فاکتور پورسل، توزیع میدان و تفکیک نرخ‌های تابشی و غیرتابشی کاملاً گریزناپذیر خواهد بود \cite{novotny,taflove}.

% ======================================== فصل ۳ ======================================================
\chapter{چارچوب نظری محاسبات الکترومغناطیسی و نظریه‌ی تابع گرین}

\section{رهیافت کلاسیک-کوانتومی و مدل نوسانگر دوقطبی}
بررسی دقیق و بنیادین برهم‌کنش میان یک گسیلنده‌ی کوانتومی (مانند یک اتم، مولکول فلورسنت، نقطه‌ی کوانتومی کلوئیدی یا مرکز نقص در شبکه‌ی بلوری) و یک نانوساختار پلاسمونی، در نگاه اول نیازمند حل کامل معادلات در چارچوب الکترودینامیک کوانتومی (\lr{QED}) است \cite{novotny}. در این چارچوب غایی، هم ماده (ترازهای انرژی گسیلنده) و هم میدان الکترومغناطیسی پیرامون آن باید به صورت عملگرهای کوانتیده (عملگرهای خلق و فنای فوتون) در نظر گرفته شوند \cite{novotny}. با این حال، حل معادلات \lr{QED} برای هندسه‌های پیچیده‌ی سه‌بعدی و محیط‌های اتلافی و پاشنده (مانند فلزات واقعی در بسامدهای نوری) از نظر تحلیلی و عددی تقریباً غیرممکن است \cite{novotny,koenderink}.

خوشبختانه، در متون مرجع نانواپتیک و بر اساس «اصل هم‌ارزی کلاسیک-کوانتومی»، نشان داده شده است که فرایند گسیل خودبه‌خودی و محاسبات مربوط به نرخ فروافت یک سامانه‌ی کوانتومی دوترازی در رژیم جفت‌شدگی ضعیف (\lr{Weak Coupling Regime}) با دقت بسیار بالا در چارچوب الکترودینامیک کلاسیک بررسی‌پذیر است \cite{novotny}. در این رهیافت تحلیلی، مشخص شده است که نرخ فروافت یک گسیلنده‌ی کوانتومی ($\Gamma$) دقیقاً متناسب با کل توان اتلافی ($P$) یک «نوسانگر دوقطبی الکتریکی کلاسیک» است \cite{novotny}. این هم‌ارزی، سنگ بنای تمامی شبیه‌سازی‌های عددی در فیزیک پلاسمونیک را تشکیل می‌دهد \cite{novotny,maier}.

برای توصیف ریاضی این سامانه‌ی کلاسیک، نقطه‌ی آغازین، معادلات ماکروسکوپیک ماکسول در حوزه‌ی زمان هستند \cite{novotny}. با فرض وابستگی زمانی هماهنگ به شکل $\exp(-i\omega t)$ برای تمامی میدان‌ها، مشتقات زمانی به ضریب $(-i\omega)$ تبدیل شده و معادلات چرخشی ماکسول (قانون القای فارادی و قانون آمپر-ماکسول) در حوزه‌ی بسامد به صورت زیر درمی‌آیند \cite{novotny}:
\begin{align}
\nabla\times\mathbf{E}(\mathbf{r},\omega)&=i\omega\mu_0\,\mathbf{H}(\mathbf{r},\omega)\,,\\
\nabla\times\mathbf{H}(\mathbf{r},\omega)&=-i\omega\epsilon_0\epsilon(\mathbf{r},\omega)\mathbf{E}(\mathbf{r},\omega)+\mathbf{J}(\mathbf{r},\omega)\,.
\end{align}
در این معادلات، $\mathbf{E}$ و $\mathbf{H}$ بردارهای میدان الکتریکی و مغناطیسی، $\mu_0$ تراوایی مغناطیسی خلأ، و $\epsilon(\mathbf{r},\omega)$ تابع دی‌الکتریک نسبیِ مختلط و وابسته به مکان محیط است که مرزهای دی‌الکتریک و نانوذره‌ی طلا را توصیف می‌کند \cite{novotny,maier}. گسیلنده‌ی کوانتومی به عنوان یک منبع جریان موضعی در نظر گرفته می‌شود. اگر گسیلنده را به عنوان یک دوقطبی نقطه‌ای نوسانگر با گشتاور دوقطبی ثابت $\mathbf{p}$ در موقعیت فضایی $\mathbf{r}_0$ مدل کنیم، قطبش این منبع به صورت $\mathbf{P}(\mathbf{r})=\mathbf{p}\,\delta(\mathbf{r}-\mathbf{r}_0)$ نوشته می‌شود (که در آن $\delta$ تابع دلتای دیراک است) \cite{novotny}. از آنجا که چگالی جریان، مشتق زمانی قطبش است ($\mathbf{J}=d\mathbf{P}/dt$)، چگالی جریان این دوقطبی به صورت جبری زیر وارد معادلات ماکسول می‌شود \cite{novotny}:
\begin{equation}
\mathbf{J}(\mathbf{r},\omega)=-i\omega\,\mathbf{p}\,\delta(\mathbf{r}-\mathbf{r}_0)\,.
\end{equation}
این چگالی جریان، همان منبع (\lr{Source}) تولیدکننده‌ی امواج الکترومغناطیسی است که رفتار گسیلنده را در حضور نانوآنتن تعیین می‌کند \cite{novotny}.

\section{معادله‌ی برداری هلمهولتز و نظریه‌ی تابع گرین دایادیک}
برای یافتن توزیع فضایی میدان الکتریکی تولیدشده توسط دوقطبی، باید معادلات جفت‌شده‌ی فارادی و آمپر را با یکدیگر ترکیب کنیم \cite{novotny}. با گرفتن عملگر کرل ($\nabla\times$) از طرفین قانون فارادی و جایگذاری قانون آمپر در آن، به یک معادله‌ی دیفرانسیل جزئی مرتبه‌ی دوم می‌رسیم که به «معادله‌ی موج ناهمگن» یا «معادله‌ی برداری هلمهولتز» (\lr{Vector Helmholtz Equation}) معروف است \cite{novotny}:
\begin{equation}
\nabla\times\big(\nabla\times\mathbf{E}(\mathbf{r})\big)-k^2(\mathbf{r})\mathbf{E}(\mathbf{r})=i\omega\mu_0\,\mathbf{J}(\mathbf{r})\,.
\end{equation}
در این معادله، $k(\mathbf{r})=\omega\sqrt{\mu_0\epsilon_0\epsilon(\mathbf{r})}$ عدد موج موضعی در محیط است \cite{novotny}. برای حل تحلیلی و سیستماتیک این معادله در حضور نانوساختارهای پیچیده، قدرتمندترین ابزار ریاضی در فیزیک ریاضی و الکترودینامیک کلاسیک، استفاده از «تابع گرین دایادیک» (\lr{Dyadic Green's Function}) است که با نماد $\mathbf{G}$ نشان داده می‌شود \cite{novotny}. از دیدگاه فیزیکی، تابع گرین دایادیک $\mathbf{G}(\mathbf{r},\mathbf{r}')$، بیانگر پاسخ ضربه (\lr{Impulse Response}) یا تابع انتقال محیط الکترومغناطیسی است؛ یعنی میدان الکتریکی برداری تولیدشده در موقعیت فضایی $\mathbf{r}$ که مستقیماً توسط یک دوقطبی نقطه‌ای واحد واقع در موقعیت $\mathbf{r}'$ ایجاد شده است \cite{novotny}.

از آنجا که گشتاور دوقطبی ($\mathbf{p}$) و میدان الکتریکی ($\mathbf{E}$) هر دو کمیت‌هایی برداری هستند، تابع گرین نمی‌تواند یک عدد اسکالر باشد، بلکه یک تانسور یا دایادیک ۳ در ۳ است که بردار $\mathbf{p}$ را به بردار $\mathbf{E}$ تصویر می‌کند \cite{novotny}. معادله‌ی دیفرانسیل حاکم بر خود تابع گرین به شکل زیر تعریف می‌شود \cite{novotny}:
\begin{equation}
\nabla\times\big(\nabla\times\mathbf{G}(\mathbf{r},\mathbf{r}')\big)-k^2(\mathbf{r})\mathbf{G}(\mathbf{r},\mathbf{r}')=\mathbf{I}\,\delta(\mathbf{r}-\mathbf{r}')\,,
\end{equation}
که در آن $\mathbf{I}$ تانسور همانی (\lr{Identity Tensor}) است. با داشتن این تابع، میدان الکتریکی کل در هر نقطه از فضا از طریق یک انتگرال‌گیری ساده بر روی حجم منبع به دست می‌آید \cite{novotny}:
\begin{equation}
\mathbf{E}(\mathbf{r})=i\omega\mu_0\int\mathbf{G}(\mathbf{r},\mathbf{r}')\,\mathbf{J}(\mathbf{r}')\,d^3\mathbf{r}'\,.
\end{equation}
با جایگذاری چگالی جریان دوقطبی نقطه‌ای ($\mathbf{J}=-i\omega\mathbf{p}\,\delta(\mathbf{r}-\mathbf{r}_0)$) در این انتگرال، رابطه‌ی بنیادین میان میدان الکتریکی موضعی و تابع گرین به دست می‌آید \cite{novotny}:
\begin{equation}
\mathbf{E}(\mathbf{r})=\omega^2\mu_0\,\mathbf{G}(\mathbf{r},\mathbf{r}_0)\,\mathbf{p}\,.
\end{equation}
درک عمیق فیزیک نهفته در پاسخ پلاسمونی، مستلزم بررسی مؤلفه‌های تابع گرین در فضای آزاد ($\mathbf{G}_0$) است \cite{novotny}. بسط ریاضی $\mathbf{G}_0$ نشان می‌دهد که میدان تولیدشده توسط دوقطبی از سه جزء اصلی تشکیل یافته است \cite{novotny}:
\begin{enumerate}
\item \textbf{میدان دور (\lr{Far-field}):} متناسب با $1/r$ که نشان‌دهنده‌ی امواج عرضیِ در حال انتشار به فضای بی‌نهایت است \cite{novotny}.
\item \textbf{میدان میانی (\lr{Intermediate-field}):} متناسب با $1/r^2$ که نقش گذار را ایفا می‌کند \cite{novotny}.
\item \textbf{میدان نزدیک (\lr{Near-field}):} متناسب با $1/r^3$ که ماهیتی شبه‌الکترواستاتیک و کاملاً میرا (\lr{Evanescent}) دارد \cite{novotny,maier}.
\end{enumerate}
در رژیم میدان نزدیک، جایی که فاصله‌ی مشاهده تا دوقطبی بسیار کوچک‌تر از طول موج است ($kr\ll1$)، جمله‌ی مکعبی ($1/r^3$) به طور مطلق بر تابع گرین غالب می‌گردد \cite{novotny}. این جمله، بیانگر انرژی به شدت محبوس‌شده و غیرانتشاری است. غلبه‌ی همین وابستگی $1/r^3$ است که حساسیت فوق‌العاده شدید سامانه به فواصل چند نانومتری از سطح نانوذره‌ی طلا را توجیه می‌کند \cite{novotny,maier}.

مهم‌ترین ویژگی کاربردی تابع گرین در مدل‌سازی سامانه‌های پلاسمونی، قابلیت تجزیه‌ی خطی آن است. بر اساس اصل برهم‌نهی، تابع گرین کل سامانه را می‌توان به دو مؤلفه‌ی مجزا تفکیک نمود \cite{novotny}:
\begin{equation}
\mathbf{G}(\mathbf{r},\mathbf{r}')=\mathbf{G}_0(\mathbf{r},\mathbf{r}')+\mathbf{G}_{\sub{scatt}}(\mathbf{r},\mathbf{r}')\,.
\end{equation}
جمله‌ی اول ($\mathbf{G}_0$) بیانگر پاسخ در فضای آزاد (خلأ) است. اما جمله‌ی دوم ($\mathbf{G}_{\sub{scatt}}$) که به «تابع گرین پراکنده» (\lr{Scattered Green's Function}) معروف است، دربردارنده‌ی تمامی تأثیرات ناشی از محیط ساختاریافته (نانوذره‌ی طلا) است \cite{novotny}. از نظر فیزیکی، $\mathbf{G}_{\sub{scatt}}$ نشان‌دهنده‌ی میدانی است که از دوقطبی ساطع شده، به مرزهای دی‌الکتریک نانوذره برخورد کرده، قطبش‌های موضعی ایجاد نموده و پس از بازتاب و پراکندگی، مجدداً به محیط بازگشته است \cite{novotny}. ارزیابی این جمله‌ی پراکنده در محل قرارگیری خود منبع ($\mathbf{r}=\mathbf{r}_0$)، تعیین‌کننده‌ی اثر پس‌کنش فیزیکی (\lr{Back-action}) و نیروی عکس‌العمل محیط بر روی خود گسیلنده می‌باشد \cite{novotny}. برای نانوکره، $\mathbf{G}_{\sub{scatt}}$ به صورت تحلیلی و با بسط چندقطبی قابل محاسبه است \cite{ruppin,mertens}؛ اما برای نانومیله‌ی کلاهک‌دار باید به صورت عددی به دست آید.

\section{چگالی حالت‌های موضعی (\lr{LDOS}) و قانون طلایی فرمی}
در فیزیک اتمی سنتی اوایل قرن بیستم، همواره تصور بر این بود که نرخ گسیل خودبه‌خودی (\lr{Spontaneous Emission Rate}) و در نتیجه طول عمر یک اتم برانگیخته، پارامتری مطلقاً ذاتی است \cite{novotny}. با این حال، بر اساس مبانی مکانیک کوانتومی وابسته به زمان و منطبق بر «قانون طلایی فرمی» (\lr{Fermi's Golden Rule})، امروزه به خوبی اثبات شده است که نرخ گذار یک سامانه‌ی کوانتومی ($\Gamma$) از حالت برانگیخته ($|i\rangle$) به حالت پایه ($|f\rangle$)، نه تنها به ویژگی‌های درونی گسیلنده (عنصر ماتریسی گذار دوقطبی)، بلکه به شدت به محیط الکترومغناطیسی اطراف بستگی دارد \cite{novotny}:
\begin{equation}
\Gamma=\frac{2\pi}{\hbar}\,\big|\langle f|\hat{H}'|i\rangle\big|^2\,\rho(\mathbf{r}_0,\omega)\,.
\end{equation}
در این رابطه، $\hat{H}'$ هامیلتونی برهم‌کنش میان اتم و میدان است و $\rho(\mathbf{r}_0,\omega)$ نمایانگر «چگالی حالت‌های موضعی» (\lr{Local Density of States - LDOS}) فوتونی است \cite{novotny}. به بیان فیزیکی، \lr{LDOS} نشان‌دهنده‌ی تعداد حالت‌های الکترومغناطیسی مجاز در دسترس در یک مختصات فضایی مشخص ($\mathbf{r}_0$) و در یک بسامد معین ($\omega$) است که فوتون گسیل‌شده می‌تواند آن‌ها را اشغال نماید \cite{novotny,maier}. در فضای آزاد (خلأ)، این چگالی حالت‌ها به صورت $\rho_0(\omega)=\omega^2/(\pi^2c^3)$ محاسبه می‌شود که منجر به نرخ فروافت ثابتِ اتم در خلأ می‌گردد \cite{novotny}.

ارتباط بنیادین میان فیزیک کوانتوم و الکترودینامیک کلاسیک دقیقاً در همین نقطه شکل می‌گیرد. در نظریه‌ی نانواپتیک به طور دقیق و تحلیلی اثبات شده است که \lr{LDOS}، از نظر ریاضی مستقیماً و بدون هیچ تقریبی با بخش موهومی تابع گرین دایادیک کلاسیک در ارتباط است \cite{novotny}. برای یک دوقطبی با بردار جهت‌گیری واحد $\mathbf{n}_p$، «چگالی حالت‌های موضعی جزئی» (\lr{Partial LDOS}) در راستای دوقطبی به صورت جبری زیر استخراج می‌گردد \cite{novotny}:
\begin{equation}
\rho_p(\mathbf{r}_0,\omega)=\frac{6\omega}{\pi c^2}\Big[\mathbf{n}_p\cdot\Imag\{\mathbf{G}(\mathbf{r}_0,\mathbf{r}_0,\omega)\}\cdot\mathbf{n}_p\Big]\,.
\end{equation}
این معادله به روشنی نشان می‌دهد که بخش موهومی تانسور گرین، زمانی که در محل استقرار خود منبع ارزیابی شود، مستقیماً شمارنده‌ی تعداد حالت‌های الکترومغناطیسی محیط است \cite{novotny}. هنگامی که یک نانوساختار مانند نانومیله یا نانوکره‌ی طلا در مجاورت گسیلنده قرار می‌گیرد، با تولید تابع گرین پراکنده ($\mathbf{G}_{\sub{scatt}}$)، چگالی حالت‌های موضعی را به شدت دچار اعوجاج می‌کند \cite{novotny,maier}. اگر هندسه و ابعاد نانوذره به گونه‌ای تنظیم شود که بسامد تشدید پلاسمون سطحی (\lr{LSPR}) با بسامد گسیلنده مطابقت داشته باشد، نانوذره تعداد بی‌شماری حالت الکترومغناطیسی جدید (مدهای پلاسمونی) برای دوقطبی فراهم می‌آورد \cite{maier,novotny}. این فوران حالت‌های در دسترس، مسیرهای جدیدی برای تخلیه‌ی انرژی دوقطبی ایجاد کرده و موجب تخلیه‌ی بسیار سریع‌تر انرژی آن نسبت به خلأ می‌گردد \cite{novotny,maier}.

\section{اثر پورسل، تفکیک نرخ‌های توان و بازده کوانتومی}
میزان تقویت نرخ گسیل خودبه‌خودی توسط محیط ساختاریافته، توسط پارامتر بدون‌بُعدی به نام «فاکتور پورسل» (\lr{Purcell Factor}، $F_P$) اندازه‌گیری می‌شود \cite{purcell,novotny}. بر اساس هم‌ارزی کلاسیک-کوانتومی که پیش‌تر تشریح شد، فاکتور پورسل معادل است با نسبت نرخ فروافت کل در حضور نانوساختار ($\Gamma_{\sub{tot}}$) به نرخ فروافت در فضای آزاد ($\Gamma_0$). این نسبت کوانتومی، دقیقاً برابر با نسبت «توان الکترومغناطیسی اتلافی کل» دوقطبی در حضور نانوآنتن ($P_{\sub{tot}}$) به توان اتلافی در خلأ ($P_0$) است \cite{novotny}.

بر اساس قوانین پایه‌ی الکترودینامیک کلاسیک، کل توان اتلافی توسط یک چگالی جریان، برابر است با کار انجام‌شده توسط میدان الکتریکی موضعی بر روی خود جریان برای حفظ نوسانات آن \cite{novotny}. این توان از انتگرال $P=-\tfrac{1}{2}\int\Real\{\mathbf{J}^*\cdot\mathbf{E}\}\,d^3\mathbf{r}$ محاسبه می‌شود. از آنجا که جریانِ دوقطبی توسط یک تابع دلتای دیراک توصیف شده است، انتگرال‌گیری ساده شده و توان کل به شکل زیر به دست می‌آید \cite{novotny}:
\begin{equation}
P_{\sub{tot}}=\frac{\omega}{2}\,\Imag\{\mathbf{p}^*\cdot\mathbf{E}(\mathbf{r}_0)\}\,.
\end{equation}
با تجزیه‌ی میدان الکتریکی کل به میدان خلأ ($\mathbf{E}_0$) و میدان پراکنده‌شده از نانوساختار ($\mathbf{E}_{\sub{scatt}}$)، می‌توان توان کل را نیز به $P_0$ و $P_{\sub{scatt}}$ تجزیه کرد. در نتیجه، فاکتور پورسل به صورت تحلیلی به شکل زیر بسط داده می‌شود \cite{novotny}:
\begin{equation}
F_P=\frac{P_{\sub{tot}}}{P_0}=1+\frac{\Imag\{\mathbf{p}^*\cdot\mathbf{E}_{\sub{scatt}}(\mathbf{r}_0)\}}{\Imag\{\mathbf{p}^*\cdot\mathbf{E}_0(\mathbf{r}_0)\}}\,.
\end{equation}
این معادله به وضوح نشان می‌دهد که تمامی تقویت نرخ فروافت (اثر پورسل)، مستقیماً برخاسته از میدان پراکنده‌ی پس‌کنش (\lr{Back-action}) از سوی نانوساختار است \cite{novotny}.

با این وجود، در بررسی نانوساختارهای پلاسمونی مانند نانومیله‌ها و نانوکره‌های طلا، یک چالش اساسی فیزیکی بروز می‌کند: افزایش چشمگیر فاکتور پورسل به هیچ وجه ضامن افزایش تولید فوتون‌های مفید در فضا نیست \cite{koenderink,novotny}. در مجاورت فلزاتِ واقعی، به دلیل وجود بخش موهومی در تابع دی‌الکتریک طلا ($\Imag(\epsilon)>0$)، امواج الکترومغناطیسی میدان نزدیک به شدت با شبکه‌ی فلز برهم‌کنش داشته و مستهلک می‌شوند \cite{maier,novotny}. بنابراین، انرژی خارج‌شده از دوقطبی ($P_{\sub{tot}}$) پس از برخورد به نانوذره، به ناچار به دو کانال کاملاً متمایز و رقیب هدایت می‌شود \cite{novotny,anger}:
\begin{equation}
P_{\sub{tot}}=P_{\sub{rad}}+P_{\sub{nonrad}}\,.
\end{equation}
\begin{enumerate}
\item \textbf{توان تابشی ($P_{\sub{rad}}$):} این مؤلفه نمایانگر انرژی الکترومغناطیسی است که با موفقیت از سامانه‌ی جفت‌شده (دوقطبی + نانوذره) گریخته و به شکل امواج عرضی فوتونی به فضای دوردست (\lr{Far-field}) تابش می‌شود \cite{novotny}. از نظر ریاضی، این توان با انتگرال‌گیری از مؤلفه‌ی عمودی بردار پوینتینگ ($\mathbf{S}=\tfrac{1}{2}\Real\{\mathbf{E}\times\mathbf{H}^*\}$) بر روی یک سطح بسته‌ی بسیار بزرگ که کل سامانه را احاطه کرده است، محاسبه می‌گردد \cite{novotny}.
\item \textbf{توان غیرتابشی ($P_{\sub{nonrad}}$):} این مؤلفه نشان‌دهنده‌ی انرژی جذب‌شده در داخل حجم فلزی نانوذره است که مستقیماً به دلیل مقاومت اهمی جریان‌های القایی و اصطکاک الکترون‌ها، به فونون و حرارت تبدیل می‌شود \cite{novotny,maier}. این توان متناسب با انتگرال حجمی $\tfrac{\omega}{2}\epsilon_0\int\Imag(\epsilon)\,|\mathbf{E}|^2\,dV$ در داخل نانوذره‌ی طلا است \cite{novotny}.
\end{enumerate}
هنگامی که گسیلنده در فواصل بسیار نزدیک (کمتر از ۵ نانومتر) از سطح نانوذره‌ی طلا قرار می‌گیرد، عبارت‌های میدان نزدیک تابع گرین پراکنده (وابستگی $1/d^3$) به شدت رشد می‌کنند \cite{novotny,anger}. در این رژیم، جفت‌شدگی میدان به قدری قوی است که مؤلفه‌ی غیرتابشی بر سامانه مسلط شده و پدیده‌ای موسوم به «خاموشی غیرتابشی» (\lr{Quenching}) رخ می‌دهد؛ در این حالت با وجود بالا بودن فاکتور پورسل، سهم تابشی از کل توان ($P_{\sub{rad}}/P_{\sub{tot}}$) به شدت کاهش می‌یابد \cite{novotny,anger}.

برای ارزیابی واقعی راندمان نانوآنتن پلاسمونی در تولید چشمه‌های تک‌فوتونی، مفهوم نرخ کل به تنهایی کافی نیست و باید از معیار فیزیکیِ «بازده کوانتومی ظاهری» (\lr{Quantum Yield}، $\eta$) استفاده شود \cite{novotny,koenderink}. این معیار بدون‌بُعد، درصد تبدیل انرژیِ خام به فوتون‌های مفید را نشان می‌دهد و به صورت زیر تعریف می‌گردد \cite{novotny}:
\begin{equation}\label{eq:eta}
\eta=\frac{P_{\sub{rad}}}{P_{\sub{tot}}}=\frac{\Gamma_{\sub{rad}}}{\Gamma_{\sub{tot}}}\,.
\end{equation}
این تعریف با فرض یک گسیلنده‌ی ایده‌آل (بازده کوانتومی ذاتی برابر یک) نوشته شده است؛ بنابراین $\eta$ در واقع «بازده آنتن» است و برای یک گسیلنده‌ی واقعی باید در بازده ذاتی آن نیز ضرب شود \cite{koenderink,novotny}.

برای نانوکره، که در این پژوهش به عنوان هندسه‌ی مرجع به کار رفته است، کمیت‌های فوق به صورت تحلیلی قابل محاسبه‌اند \cite{gersten,ruppin}. برای دوقطبی‌ای که در فاصله‌ی $r$ از مرکز یک کره به شعاع $a$ و عمود بر سطح آن (در راستای شعاعی) قرار دارد، در حد شبه‌استاتیک و با در نظر گرفتن تمام مرتبه‌های چندقطبی، داریم \cite{ruppin,mertens}:
\begin{equation}\label{eq:sphere-rates}
F_P^{\perp}=1+\frac{3}{2k^3}\sum_{n=1}^{\infty}(n+1)^2\,\frac{\Imag\{\beta_n\}}{r^{2n+4}},\qquad
\frac{\Gamma_{\sub{rad}}^{\perp}}{\Gamma_0}=\Big|1+\frac{2\beta_1}{r^3}\Big|^2,\qquad
\beta_n=\frac{n(\epsilon_m-\epsilon_d)}{n\epsilon_m+(n+1)\epsilon_d}\,a^{2n+1}\,,
\end{equation}
و برای دوقطبی مماس بر سطح، ضریب $(n+1)^2$ به $n(n+1)/2$ و جمله‌ی تابشی به $|1-\beta_1/r^3|^2$ تبدیل می‌شود \cite{ruppin}. جمله‌ی $n=1$ همان پاسخ دوقطبی کره (رابطه‌ی (\ref{eq:alpha-sphere}) با $\alpha=4\pi\beta_1$) است، و جمله‌های مرتبه‌ی بالاتر، که مدهای «تاریک» و غیرتابشی کره را توصیف می‌کنند، در فواصل کوتاه غالب شده و عامل اصلی خاموشی هستند \cite{ruppin,mertens}. این روابط، معیاری دقیق و مستقل برای مقایسه‌ی نتایج عددی نانومیله با یک هندسه‌ی ساده‌تر فراهم می‌کنند.

در شبیه‌سازی‌های دقیق الکترومغناطیسی مبتنی بر روش \lr{FDTD} در فصول آتی، رسالت اصلی کدهای محاسباتی، تفکیک دقیق این دو انتگرال (سطحی برای $P_{\sub{rad}}$ و حجمی برای $P_{\sub{nonrad}}$) است تا بتوان با پویش پیوسته‌ی فواصل فضایی، نقطه‌ی بهینه‌ای را یافت که در آن نانومیله‌ی طلا به عنوان یک «آنتن نوری کارآمد» عمل کرده و بازده کوانتومی $\eta$ را در کنار یک فاکتور پورسل بالا، بیشینه سازد \cite{novotny,maier}.

% ======================================== فصل ۴ ======================================================
\chapter{الگوریتم تفاضل محدود در حوزه‌ی زمان (\lr{FDTD})}

\section{اصول روش \lr{FDTD} و گسسته‌سازی معادلات ماکسول}
بررسی تحلیلی برهم‌کنش نور و ماده، همان‌طور که در فصول پیشین به تفصیل بحث شد، تنها برای هندسه‌های بسیار متقارن و ایده‌آل مانند کره‌های همگن (از طریق نظریه‌ی می و بسط‌های چندقطبی) یا استوانه‌های بی‌نهایت امکان‌پذیر است \cite{novotny,maier,mie}. با شکسته شدن این تقارن هندسی در ساختارهای سه‌بعدی و واقعی مانند نانومیله‌های کلاهک‌دار، نانوآنتن‌های پاپیونی و سطوح زبر، حل دقیق و تحلیلی معادلات ماکسول و یافتن شکل بسته برای تانسور تابع گرین پراکنده کاملاً غیرممکن می‌گردد \cite{novotny}. در چنین شرایطی، استفاده از روش‌های حل عددی (\lr{Numerical Methods}) در الکترومغناطیس محاسباتی تنها راهکار فیزیک‌دانان و مهندسان است \cite{taflove,novotny}. حل تحلیلی نانوکره (رابطه‌ی (\ref{eq:sphere-rates})) در این میان نقش یک مرجع مستقل را ایفا می‌کند که روندهای کلی به‌دست‌آمده از شبیه‌سازی را می‌توان با آن سنجید.

در میان روش‌های عددی متنوع مانند روش المان محدود (\lr{FEM}) و روش ممان‌ها (\lr{MoM}) که عمدتاً معادلات را در حوزه‌ی بسامد حل می‌کنند، «روش تفاضل محدود در حوزه‌ی زمان» (\lr{Finite-Difference Time-Domain - FDTD}) به دلیل ماهیت ذاتی خود، جایگاه ویژه‌ای در فیزیک نانواپتیک و پلاسمونیک یافته است \cite{taflove,novotny}. روش \lr{FDTD} که یک رویکرد مستقیم برای حل معادلات دیفرانسیل جزئی وابسته به زمان ماکسول است، دو مزیت بی‌بدیل برای بررسی سامانه‌های پلاسمونی فراهم می‌آورد: نخست آنکه با استفاده از تنها یک پالس تحریکی پهن‌باند (\lr{Broadband Gaussian Pulse}) در حوزه‌ی زمان و سپس اعمال تبدیل فوریه‌ی گسسته (\lr{DFT}) بر روی پاسخ سامانه، می‌توان طیف نوری ساختار (مانند طیف جذب، سطح مقطع پراکندگی و نمودار فاکتور پورسل) را به صورت هم‌زمان در یک بازه‌ی وسیع بسامدی استخراج نمود \cite{taflove}. دوم آنکه این روش به هیچ تقارن خاصی در هندسه نیاز ندارد و مدل‌سازی ویژگی‌های پاشندگی پیچیده‌ی مواد (مانند وابستگی شدید تابع دی‌الکتریک طلا به بسامد) را بدون نیاز به وارون‌سازی ماتریس‌های بزرگ، امکان‌پذیر می‌سازد \cite{taflove,maier}.

نقطه‌ی آغازین الگوریتم \lr{FDTD}، معادلات چرخشی ماکسول در یک محیط خطی، ایزوتروپیک و غیرمغناطیسی (تراوایی $\mu=\mu_0$) است؛ در این فصل، $\epsilon$ گذردهی مطلق محیط ($\epsilon_0\epsilon_r$) را نشان می‌دهد. قانون القای فارادی و قانون آمپر-ماکسول به شکل دیفرانسیلی زیر نوشته می‌شوند \cite{taflove,novotny}:
\begin{align}
\frac{\partial\mathbf{H}}{\partial t}&=-\frac{1}{\mu_0}\nabla\times\mathbf{E}\,,\\
\frac{\partial\mathbf{E}}{\partial t}&=\frac{1}{\epsilon}\nabla\times\mathbf{H}-\frac{1}{\epsilon}\mathbf{J}\,.
\end{align}
در یک دستگاه مختصات دکارتی سه‌بعدی ($x,y,z$)، عملگر کرل (\lr{Curl}) باعث می‌شود که این دو معادله‌ی برداری، به یک دستگاه متشکل از شش معادله‌ی دیفرانسیل جزئی جفت‌شده تبدیل گردند \cite{yee,taflove}. به عنوان نمونه، معادله‌ی مربوط به مؤلفه‌ی $x$ از میدان الکتریکی ($E_x$) و مؤلفه‌ی $x$ از میدان مغناطیسی ($H_x$) به صورت زیر بسط داده می‌شوند \cite{yee,taflove}:
\begin{align}
\frac{\partial E_x}{\partial t}&=\frac{1}{\epsilon}\left(\frac{\partial H_z}{\partial y}-\frac{\partial H_y}{\partial z}-J_x\right)\,,\\
\frac{\partial H_x}{\partial t}&=\frac{1}{\mu_0}\left(\frac{\partial E_y}{\partial z}-\frac{\partial E_z}{\partial y}\right)\,.
\end{align}
مشاهده می‌شود که تغییرات فضایی میدان مغناطیسی (گرادیان‌های $H_z$ و $H_y$) تعیین‌کننده‌ی نرخ تغییرات زمانی میدان الکتریکی ($E_x$) هستند و بالعکس. این جفت‌شدگی فضا-زمان، اساس طراحی شبکه‌ی محاسباتی \lr{FDTD} را تشکیل می‌دهد \cite{yee,taflove}.

\section{شبکه‌ی یی (\lr{Yee Grid}) و الگوریتم جهشی (\lr{Leapfrog})}
قلب تپنده و بنیان ریاضی روش \lr{FDTD}، الگوریتمی است که در سال ۱۹۶۶ توسط کین یی (\lr{Kane Yee}) معرفی گردید \cite{yee}. ایده‌ی درخشان «الگوریتم یی»، گسسته‌سازی (\lr{Discretization}) هوشمندانه‌ی این شش معادله‌ی جفت‌شده در یک شبکه‌ی فضایی-زمانی است \cite{yee,taflove}. یی به جای آنکه تمامی مؤلفه‌های میدان ($E_x,E_y,E_z,H_x,H_y,H_z$) را در یک گره (\lr{Node}) واحد از شبکه قرار دهد، فضا را به مکعب‌های کوچکی به نام «سلول‌های یی» (\lr{Yee Cells}) با ابعاد $\Delta x,\Delta y,\Delta z$ تقسیم نمود و مؤلفه‌ها را به صورت متقاطع در این سلول توزیع کرد \cite{yee}.

نوآوری اصلی این الگوریتم در چیدمان هندسی میدان‌هاست؛ مؤلفه‌های میدان الکتریکی دقیقاً بر روی وسط یال‌های مکعب، و مؤلفه‌های میدان مغناطیسی دقیقاً در مرکز وجه‌های مکعب قرار می‌گیرند \cite{yee,taflove}. این معماری شبکه‌ای باعث می‌شود که هر مؤلفه‌ی میدان $\mathbf{E}$ توسط حلقه‌ای متشکل از چهار مؤلفه‌ی میدان $\mathbf{H}$ احاطه شود و متقابلاً هر مؤلفه‌ی $\mathbf{H}$ توسط چهار مؤلفه‌ی $\mathbf{E}$ محاصره گردد. از نظر ریاضی، این ساختار درهم‌تنیده امکان می‌دهد که مشتقات جزئی فضایی در معادلات ماکسول، با استفاده از تقریب «تفاضل مرکزی» (\lr{Central Finite Difference}) با دقتی از مرتبه‌ی دوم ($O(\Delta^2)$) نسبت به اندازه‌ی شبکه تخمین زده شوند \cite{yee,taflove}.

به عنوان مثال، تقریب تفاضل مرکزی برای مشتق یک تابع پیوسته‌ی $f(x)$ در نقطه‌ی $x_0$ با گام $\Delta x$ به صورت $\big(f(x_0+\Delta x/2)-f(x_0-\Delta x/2)\big)/\Delta x$ نوشته می‌شود. شبکه‌ی یی این نیم‌گام‌های فضایی ($\Delta x/2$) را به صورت ساختاری در خود جای داده است \cite{yee}.

علاوه بر جابه‌جایی فضایی، نبوغ کین یی در بُعد زمان نیز اعمال شد. میدان‌های $\mathbf{E}$ و $\mathbf{H}$ در زمان نیز به اندازه‌ی نیم‌گام زمانی ($\Delta t/2$) نسبت به یکدیگر تغییر مکان داده می‌شوند \cite{yee,taflove}. یعنی اگر میدان‌های الکتریکی در زمان‌های صحیح ($t=n\Delta t$) محاسبه شوند، میدان‌های مغناطیسی در زمان‌های نیم‌صحیح ($t=(n+1/2)\Delta t$) ارزیابی می‌گردند. این شیوه، که در آنالیز عددی به «الگوریتم جهشی» (\lr{Leapfrog Integration}) معروف است، سبب می‌شود که میدان‌های الکتریکی و مغناطیسی به طور متناوب و مانند یک زنجیره‌ی علت و معلولی در طول زمان به پیش بروند \cite{taflove}.

با اعمال این گسسته‌سازی‌ها بر روی معادله‌ی دیفرانسیل جزئی $E_x$ (که در بخش قبل بیان شد)، فرمول گسسته‌شده و صریح زیر برای به‌روزرسانی مؤلفه‌ی $E_x$ در گام زمانی آینده ($n+1$) بر اساس مقادیر فعلی آن در گام ($n$) و مقادیر پیرامونی میدان مغناطیسی در گام میانی ($n+1/2$) به دست می‌آید \cite{taflove}:
\begin{multline}
E_x\big|^{n+1}_{i,j+\frac12,k+\frac12}=E_x\big|^{n}_{i,j+\frac12,k+\frac12}+\frac{\Delta t}{\epsilon_{i,j+\frac12,k+\frac12}}\\
\times\left[\frac{H_z\big|^{n+\frac12}_{i,j+1,k+\frac12}-H_z\big|^{n+\frac12}_{i,j,k+\frac12}}{\Delta y}-\frac{H_y\big|^{n+\frac12}_{i,j+\frac12,k+1}-H_y\big|^{n+\frac12}_{i,j+\frac12,k}}{\Delta z}-J_x\big|^{n+\frac12}_{i,j+\frac12,k+\frac12}\right]
\end{multline}
در این معادله‌ی جبری، اندیس‌های $i,j,k$ نشان‌دهنده‌ی مختصات گره‌های شبکه‌ی فضایی هستند. این معادله نشان می‌دهد که حل‌کننده \lr{FDTD} برای پیش‌بینی آینده‌ی سامانه نیازی به وارون‌سازی ماتریس ندارد \cite{taflove}. با شروع از شرایط اولیه ($t=0$) که در آن تمام میدان‌ها صفر فرض می‌شوند، یک منبع جریانی ($J$) سامانه را تحریک کرده و رایانه با یک حلقه‌ی تکرار (\lr{For Loop}) ساده، سلول به سلول فضای محاسباتی را جاروب نموده و میدان‌ها را در زمان به جلو می‌راند \cite{taflove}.

\section{پایداری عددی و شرط کورانت-فردریشز-لوی (\lr{CFL})}
با وجود توانمندی بالا و سادگی پیاده‌سازی الگوریتم صریح یی، این روش از نظر ریاضی «پایدار مشروط» (\lr{Conditionally Stable}) است \cite{taflove}. به این معنا که انتخاب اندازه‌ی گام‌های مکانی ($\Delta x,\Delta y,\Delta z$) و گام زمانی ($\Delta t$) نمی‌تواند به صورت کاملاً مستقل و دلخواه صورت پذیرد. اگر پارامترهای گسسته‌سازی به درستی تنظیم نشوند، خطاهای گردکردن عددی (\lr{Round-off errors}) در هر گام زمانی انباشته شده و به صورت نمایی رشد می‌کنند که منجر به انفجار محاسباتی و رسیدن مقادیر میدان‌ها به سمت بی‌نهایت می‌گردد (پدیده‌ی واگرایی یا \lr{Divergence}) \cite{taflove}.

از دیدگاه فیزیکی، امواج الکترومغناطیسی حداکثر با سرعت نور در خلأ ($c$) منتشر می‌شوند. برای آنکه الگوریتم جهشی یی بتواند انتقال اطلاعات (انتشار فاز موج) از یک سلول محاسباتی به سلول مجاور را به درستی و بر اساس اصل علیت ردگیری کند، ضروری است که مسافت فیزیکی طی‌شده توسط موج در طول یک گام زمانی مجزا ($c\Delta t$)، هرگز از ابعاد هندسی کوچک‌ترین سلول محاسباتی تجاوز نکند \cite{taflove}. اگر موج در بازه‌ی $\Delta t$ از بیش از یک سلول عبور کند، الگوریتم عملاً از موج جا مانده و معادلات تفاضلی ارتباط منطقی خود را با معادلات دیفرانسیل پیوسته از دست می‌دهند.

این محدودیت بنیادین در آنالیز عددی معادلات هذلولوی (\lr{Hyperbolic PDEs})، تحت عنوان «شرط پایداری کورانت-فردریشز-لوی» (\lr{Courant-Friedrichs-Lewy Condition}) یا به اختصار شرط کورانت (\lr{CFL}) شناخته می‌شود \cite{taflove}. با انجام تحلیل پایداری فون‌نویمان (\lr{Von Neumann Stability Analysis}) بر روی معادلات گسسته‌ی یی، شرط پایداری برای یک شبکه‌ی سه‌بعدی دکارتی به صورت نامعادله‌ی اکید زیر اثبات می‌گردد \cite{taflove}:
\begin{equation}
\Delta t\le\frac{1}{c\sqrt{\dfrac{1}{(\Delta x)^2}+\dfrac{1}{(\Delta y)^2}+\dfrac{1}{(\Delta z)^2}}}\,.
\end{equation}
این رابطه تأثیرات مهمی در شبیه‌سازی‌های پلاسمونیک دارد \cite{maier,taflove}. برای یک موج نوری (مثلاً با طول موج ۶۰۰ نانومتر)، به طور معمول شبکه‌بندی فضایی با ابعاد $\Delta x=10$ نانومتر برای حل معادلات کافی است. اما در حوزه‌ی نانوفوتونیک و در بررسی اثر پورسل، به دلیل تمرکز شدید میدان در نوک نانومیله‌ها و نیاز حیاتی به ردیابی افت سریع میدان نزدیک در فواصل ۳ تا ۲۰ نانومتر از سطح فلز، ناچاریم از یک مش‌بندی موضعی بسیار ظریف در حد نانومتر (مثلاً $\Delta x=\Delta y=\Delta z=1$ نانومتر) استفاده کنیم \cite{novotny,taflove}.

با جایگذاری $\Delta x=1$ نانومتر در رابطه‌ی \lr{CFL}، بیشینه‌ی گام زمانی مجاز برای پایدار ماندن الگوریتم به حدود $\Delta t\approx1.9\times10^{-18}$ ثانیه (۱/۹ آتوثانیه) کاهش می‌یابد. این بدان معناست که برای شبیه‌سازی تنها چند صد فمتوثانیه از دینامیک فروافت دوقطبی، رایانه باید صدها هزار بار حلقه‌ی محاسباتی سه‌بعدی را تکرار نماید. این مصالحه میان «دقت فضایی» و «هزینه‌ی محاسباتی زمانی»، بزرگ‌ترین چالش شبیه‌سازی \lr{FDTD} در پلاسمونیک است \cite{taflove,maier}.

\section{شرایط مرزی در شبیه‌سازی پلاسمونیک و لایه‌های \lr{PML}}
یکی دیگر از چالش‌های بنیادین در حل عددی معادلات دیفرانسیل جزئی برای مسائل تابشی، محدود بودن فضای در دسترس رایانه (حافظه‌ی \lr{RAM}) است \cite{taflove}. یک دوقطبی نقطه‌ای که در مجاورت نانومیله نوسان می‌کند، امواج الکترومغناطیسی (فوتون‌ها) را به فضای بیکران و نامتناهی ساطع می‌نماید. اما در شبیه‌سازی \lr{FDTD}، ما مجبوریم ناحیه‌ی محاسباتی را در یک حجم سه‌بعدی محدود (\lr{Simulation Region}) قطع و محصور کنیم. اگر در مرزهای این فضای برش‌خورده از شرایط مرزی مناسبی استفاده نشود، امواجی که به دیواره‌های شبکه می‌رسند، به صورت غیرفیزیکی بازتاب شده و مجدداً به داخل محیط محاسباتی بازمی‌گردند و با میدان اصلی تداخل می‌کنند؛ پدیده‌ای که به کلی نتایج محاسبه‌ی فاکتور پورسل، الگوی تابشی و چگالی حالت‌های موضعی را مخدوش می‌سازد \cite{taflove,novotny}.

برای رفع این معضل در مسائل پراکندگی فضای باز (\lr{Open Boundaries})، قدرتمندترین روش موجود که امروزه به استاندارد طلایی \lr{FDTD} تبدیل شده است، استفاده از «لایه‌های کاملاً تطبیق‌یافته» (\lr{Perfectly Matched Layers - PML}) است که در سال ۱۹۹۴ توسط ژان‌پیر برنجر (\lr{J. P. Berenger}) ابداع گردید \cite{berenger}. برخلاف شرایط مرزی ساده‌ی جذب‌کننده‌ی تحلیلی (\lr{ABC})، لایه‌ی \lr{PML} یک شرط مرزی صرف نیست، بلکه یک «محیط جاذب مصنوعی و ریاضی» با ضخامت چند سلول محاسباتی است که در دورتادور مرزهای فضای اصلی قرار داده می‌شود \cite{berenger,taflove}.

ویژگی بی‌نظیر \lr{PML} در فرمول‌بندی ریاضی آن نهفته است که بر پایه‌ی روش «کشش مختصات مختلط» (\lr{Complex Coordinate Stretching}) بنا شده است \cite{taflove}. در این روش، مختصات فضایی در داخل ناحیه‌ی \lr{PML} به مقادیر مختلط تبدیل می‌شوند (به صورت $x\to x+i\sigma_x/\omega$). این تبدیل ریاضی باعث می‌شود که امپدانس موجی (\lr{Wave Impedance}) لایه‌ی \lr{PML} در حد پیوسته (پیش از گسسته‌سازی) در تمامی بسامدها و زوایای تابش، با امپدانس محیط محاسباتی مجاور (مانند هوا یا شیشه) تطبیق یابد \cite{berenger,taflove}.

به واسطه‌ی این تطبیق کامل امپدانس، امواج ساطع‌شده از سامانه‌ی دوقطبی-نانومیله، به محض رسیدن به لایه‌ی \lr{PML}، تقریباً بدون بازتاب (بازتاب عددی بسیار ناچیز) وارد این لایه‌ی مصنوعی می‌شوند. پس از ورود، رسانش الکتریکی و مغناطیسی مصنوعی تعریف‌شده در \lr{PML} ($\sigma_x,\sigma_x^*$)، انرژی امواج را به صورت نمایی و با ضریب تضعیف بسیار بالا، در طول مسیر انتشار مستهلک کرده و پیش از رسیدن به دیواره‌ی صلب نهایی (مرز \lr{PEC})، دامنه‌ی موج را به صفر می‌رسانند \cite{berenger,taflove}. در شبیه‌سازی‌های این پایان‌نامه، مرزهای محاسباتی در هر ۶ وجه ساختار به عنوان \lr{PML} استاندارد تنظیم شده‌اند تا اطمینان حاصل گردد که نرخ تابشی ($\Gamma_{\sub{rad}}$) استخراج‌شده برای دوقطبی، به درستی نمایانگر انتشار فوتون‌ها به فضای بی‌نهایت است و هیچ‌گونه تداخل ناشی از بازتاب عددی در سامانه وجود ندارد.

\section{مدل‌سازی پاشندگی مواد (\lr{Material Dispersion}) در \lr{FDTD}}
همان‌طور که در فصل دوم بحث شد، وارد کردن فیزیک فلزات نجیب مانند طلا در شبیه‌سازی‌های پلاسمونیک امری حیاتی است \cite{maier}. گذردهی الکتریکی (\lr{Permittivity}) طلا یک عدد ثابت نیست، بلکه در رژیم نوری پاشندگی شدیدی دارد و به شدت وابسته به بسامد است: $\epsilon(\omega)=\epsilon'(\omega)+i\epsilon''(\omega)$ \cite{maier,jc}.

الگوریتم استاندارد یی در حوزه‌ی زمان بر اساس رابطه‌ی ساختاری ساده‌ی $\mathbf{D}=\epsilon\mathbf{E}$ (با $\epsilon$ ثابت) بنا شده است. اما در محیط‌های پاشنده، تابع دی‌الکتریک به زمان وابسته است و این ضرب در حوزه‌ی بسامد، به واسطه‌ی قضیه‌ی کانولوشن، به یک «انتگرال کانولوشن» در حوزه‌ی زمان تبدیل می‌گردد \cite{taflove}:
\begin{equation}
\mathbf{D}(t)=\epsilon_0\epsilon_\infty\mathbf{E}(t)+\epsilon_0\int_0^t\chi(\tau)\,\mathbf{E}(t-\tau)\,d\tau\,,
\end{equation}
که در آن $\chi(\tau)$ تبدیل فوریه‌ی وارون پذیرفتاری الکتریکی (\lr{Susceptibility}) است. حل مستقیم این انتگرال کانولوشن در هر گام زمانی شبکه‌ی \lr{FDTD}، نیازمند ذخیره‌سازی تمامی تاریخچه‌ی مقادیر میدان الکتریکی از زمان صفر تا زمان $t$ است که این امر به حافظه‌ی بسیار زیاد (\lr{Memory overflow}) و افزایش شدید زمان محاسبات منجر می‌شود \cite{taflove}.

برای غلبه بر این محدودیت بنیادین، از روش قدرتمند «معادلات دیفرانسیل کمکی» (\lr{Auxiliary Differential Equation - ADE}) استفاده می‌شود \cite{taflove}. در این روش، به جای حل مستقیم انتگرال کانولوشن، قطبش پاشنده‌ی محیط ($\mathbf{P}_{\sub{disp}}$) به عنوان یک متغیر حالت جدید در نظر گرفته می‌شود. بر اساس مدل‌های پیوسته‌ی «درود-لورنتز» (\lr{Drude-Lorentz}) یا مدل‌های چندقطبی (\lr{Multi-pole})، تابع دی‌الکتریک فلز به مجموعه‌ای از قطب‌های تشدیدی تجزیه می‌شود \cite{taflove,maier}:
\begin{equation}
\epsilon(\omega)=\epsilon_\infty+\sum_{m=1}^{M}\frac{\Delta\epsilon_m\,\omega_{0,m}^2}{\omega_{0,m}^2-\omega^2-i\omega\Gamma_m}\,.
\end{equation}
هر یک از این عبارت‌های کسری در حوزه‌ی زمان، متناظر با یک معادله‌ی دیفرانسیل معمولی (\lr{ODE}) درجه دوم برای متغیر کمکی $\mathbf{P}_m$ است. روش \lr{ADE} این معادلات کمکی را در کنار معادلات اصلی ماکسول قرار داده و در هر گام زمانی الگوریتم یی، مقادیر $\mathbf{P}_m$ را هم‌زمان با میدان‌های $\mathbf{E}$ و $\mathbf{H}$ به‌روزرسانی می‌کند \cite{taflove}. این ترفند ریاضی، نیاز به ذخیره‌ی تاریخچه‌ی زمانی را از بین برده و امکان شبیه‌سازی مواد پاشنده را با هزینه‌ی محاسباتی معادل با مواد غیرپاشنده فراهم می‌سازد \cite{taflove}. در نرم‌افزار \lr{Lumerical FDTD} که در این پایان‌نامه به کار رفته است، داده‌های تجربی جانسون-کریستی \cite{jc} برای طلای توده‌ای با یک مدل چندضریبی (\lr{Multi-Coefficient Model - MCM}) در بازه‌ی طیفی شبیه‌سازی برازش (\lr{Fit}) می‌شوند و این مدل، که از همان خانواده‌ی قطب‌های تشدیدی است، به همین روش در حوزه‌ی زمان پیاده‌سازی می‌شود تا گذارهای بین‌نواری و اتلاف اهمی طلا با دقت در محاسبات وارد شوند \cite{jc,taflove}.

\section{راه‌اندازی دوقطبی کلاسیک و استخراج فاکتور پورسل}
گام نهایی در تنظیم معماری محاسباتی \lr{FDTD} برای مدل‌سازی الکترودینامیک کوانتومی، تعیین دقیق منابع تحریکی (\lr{Sources}) و استخراج داده‌های خروجی از طریق مانیتورها (\lr{Monitors}) است \cite{novotny,taflove}. در این ساختار، گسیلنده‌ی کوانتومی با یک \textbf{منبع دوقطبی کلاسیک} (\lr{Dipole Source}) جایگزین می‌شود \cite{novotny}. این منبع به عنوان یک چگالی جریان نقطه‌ای با یک توزیع گاوسی بسیار باریک در حوزه‌ی زمان به سامانه اعمال می‌شود. به واسطه‌ی تبدیل فوریه، یک پالس باریک در زمان، طیف بسامدی پهن‌باندی را در حوزه‌ی بسامد تولید می‌کند که به حل‌کننده اجازه می‌دهد پاسخ سامانه (طیف پورسل) را در بازه‌ی طول‌موج‌های ۵۰۰ تا ۹۰۰ نانومتر در یک بار اجرای واحد (\lr{Single simulation run}) محاسبه نماید \cite{taflove}.

برای استخراج کمّی و دقیق «اثر پورسل» و تفکیک نرخ‌های توان تابشی و غیرتابشی ($\Gamma_{\sub{rad}}$ و $\Gamma_{\sub{nonrad}}$)، از روش «گروه‌های تحلیل انتقال» (\lr{Transmission Analysis Groups}) بهره گرفته شده است \cite{novotny}. توان الکترومغناطیسی عبوری از یک سطح بسته ($S$) در حوزه‌ی بسامد، به صورت انتگرال‌گیری فضایی از مؤلفه‌ی عمودی بردار پوینتینگ محاسبه می‌شود \cite{novotny}:
\begin{equation}
P(\omega)=\frac{1}{2}\oint_S\Real\{\mathbf{E}(\mathbf{r},\omega)\times\mathbf{H}^*(\mathbf{r},\omega)\}\cdot d\mathbf{a}\,.
\end{equation}
در نرم‌افزار محاسباتی لومریکال، به منظور تفکیک جریان‌های توان، دو مجموعه از جعبه‌های مانیتورینگ فضایی (هر جعبه متشکل از ۶ مانیتور صفحه‌ای در وجوه $+x,-x,+y,-y,+z,-z$) تعریف گردیده‌اند:
\begin{enumerate}
\item \textbf{جعبه‌ی مانیتورینگ منبع (\lr{Purcell Box}):} یک مکعب بسیار کوچک مانیتورینگ (با ابعادی تنها چند برابر اندازه‌ی شبکه و کوچک‌تر از فاصله‌ی گسیلنده تا فلز) که دقیقاً دورتادور خود منبع دوقطبی کشیده می‌شود، به طوری که هیچ بخشی از فلز درون آن قرار نگیرد. این مانیتور با انتگرال‌گیری از بردار پوینتینگ در همسایگی بی‌واسطه‌ی منبع، «کل توان الکترومغناطیسی خارج‌شده از دوقطبی» ($P_{\sub{tot}}$) را در حضور نانومیله و تحت تأثیر نیروی پس‌کنش تابع گرین پراکنده اندازه می‌گیرد \cite{novotny}. نسبت این توان استخراج‌شده به توان تحلیلی دوقطبی در خلأ ($P_0$)، «فاکتور پورسل» کل ساختار ($F_P=P_{\sub{tot}}/P_0$) را در هر بسامد به دست می‌دهد \cite{novotny,purcell}.
\item \textbf{جعبه‌ی مانیتورینگ تابشی (\lr{Radiative Box}):} یک مکعب بسیار بزرگ‌تر که در مجاورت لایه‌های جاذب \lr{PML} قرار گرفته و کل سامانه‌ی فیزیکی (متشکل از نانومیله‌ی طلا و دوقطبی نقطه‌ای) را در بر می‌گیرد. این مانیتور بیرونی، تنها انرژی فوتون‌هایی را می‌سنجد که از چنگال اتلاف اهمی فلز گریخته و به فضای دوردست تابش شده‌اند ($P_{\sub{rad}}$) \cite{novotny,koenderink}.
\end{enumerate}
اختلاف انرژی محاسبه‌شده میان این دو جعبه، مستقیماً نمایانگر انرژی بلعیده‌شده و تلف‌شده در نانومیله‌ی طلا (توان غیرتابشی $P_{\sub{nonrad}}$) است \cite{novotny}. با تقسیم خروجی عددی جعبه‌ی دوم (تابشی) بر جعبه‌ی اول (کل)، پارامتر حیاتی «بازده کوانتومی ظاهری» ($\eta=P_{\sub{rad}}/P_{\sub{tot}}$) در تمامی بازه‌ی بسامدی استخراج شده و راندمان حقیقی نانوآنتن پلاسمونی به عنوان یک مبدل امواج موضعی به امواج آزاد، به صورت کاملاً کمّی ارزیابی می‌گردد \cite{novotny,maier}.

% ======================================== فصل ۵ ======================================================
\chapter{نتایج شبیه‌سازی و تحلیل داده‌ها}

\section{مقدمه}
پس از تدوین چارچوب نظری و معرفی دقیق الگوریتم‌های محاسباتی در فصول پیشین، این فصل به ارائه، تشریح و تحلیل فیزیکی نتایج به دست آمده از شبیه‌سازی‌های سه‌بعدی برهم‌کنش میان یک گسیلنده‌ی کوانتومی و نانوآنتن پلاسمونی طلا اختصاص دارد \cite{novotny,maier}. در الکترودینامیک نانومقیاس، رفتار سامانه‌های جفت‌شده به شدت تحت تأثیر پارامترهای هندسی، ویژگی‌های ماده و از همه مهم‌تر، فاصله‌ی فضایی میان منبع تابش و سطح فلز قرار دارد \cite{novotny,maier}. از این رو، هدف اصلی این فصل، کالبدشکافی دقیق پدیده‌ی «اثر پورسل» در رژیم میدان نزدیک و بررسی رقابت میان «تقویت نرخ گسیل تابشی» و «خاموشی غیرتابشی» (اتلاف اهمی) است \cite{novotny,anger}. تحلیل‌های ارائه‌شده در این فصل، بدون اتکا به تقریب‌های ساده‌کننده‌ی تحلیلی و مستقیماً بر اساس حل عددی و کامل معادلات ماکسول (توسط الگوریتم \lr{3D-FDTD}) استخراج شده‌اند تا در نهایت، نقطه‌ی کار بهینه (\lr{Optimal operating point}) برای طراحی چشمه‌های تک‌فوتونی به صورت کمّی مشخص گردد \cite{novotny,taflove}. برای آنکه سهم هندسه‌ی میله‌ای از سهم خود فلز طلا جدا شود، نتایج در هر بخش با یک نانوکره‌ی طلای هم‌حجم، که مرجع کلاسیک این حوزه است، در شرایط یکسان مقایسه می‌شوند.

\section{تعریف هندسه‌ی مسئله، ویژگی‌های ماده و پارامترهای شبیه‌سازی}
گام نخست در هر شبیه‌سازی الکترومغناطیسی، تعریف دقیق فضای محاسباتی، هندسه‌ی پراکنده‌ساز و توزیع فضایی ضریب شکست است. به منظور بررسی دقیق برهم‌کنش میان گسیلنده و نانوآنتن، یک نانومیله‌ی طلا (\lr{Gold Nanorod}) با هندسه‌ی استوانه‌ای کلاهک‌دار (\lr{Hemispherically-capped cylinder}) در محیط نرم‌افزار پیاده‌سازی گردید \cite{maier,giannini}. انتخاب این هندسه به هیچ وجه تصادفی نبوده است؛ نانومیله‌های کلاهک‌دار به دلیل فقدان لبه‌های کاملاً تیز (برخلاف مکعب‌ها یا منشورها)، از ایجاد تکینگی‌های غیرفیزیکی (\lr{Singularities}) در محاسبات عددی میدان جلوگیری می‌کنند و در عین حال، انحنای مناسب در کلاهک‌ها، اثر برق‌گیر (\lr{Lightning-rod effect}) را برای تولید نقاط داغ پلاسمونی به بهترین شکل فراهم می‌آورد \cite{novotny,bharadwaj}.

ابعاد فیزیکی این نانومیله شامل طول کل ($L$) برابر با ۶۰ نانومتر و قطر مقطع عرضی ($D$) برابر با ۲۰ نانومتر در نظر گرفته شده است. این ابعاد، شعاع کلاهک‌های انتهایی را دقیقاً برابر با ۱۰ نانومتر تعیین می‌کند. نسبت طول به عرض (\lr{Aspect Ratio}) در این ساختار برابر با ۳ است \cite{maier}. این نسبت هندسی به طور خاص انتخاب شده است تا بسامد تشدید پلاسمون سطحی طولی (\lr{Longitudinal LSPR}) از ناحیه‌ی جذب شدید طلا (ناحیه‌ی سبز و آبی) دور شده و به سمت طول‌موج‌های بلندتر (انتقال به سرخ) در محدوده‌ی طیفی ۶۰۰ تا ۶۵۰ نانومتر هدایت شود \cite{maier,novotny}؛ ناحیه‌ای که برای جفت‌شدگی با بسیاری از مراکز رنگی (مانند مرکز \lr{NV} در الماس) و نقاط کوانتومی کلوئیدی مناسب است \cite{novotny}. حجم این نانومیله (یک استوانه به طول $L-2R=40$ نانومتر و دو نیم‌کره به شعاع $R=10$ نانومتر) برابر است با
\[
V=\pi R^2(L-2R)+\frac{4}{3}\pi R^3=\frac{16000}{3}\pi\approx16755\ \mathrm{nm}^3\,.\] بنابراین نانوکره‌ی مرجع، کره‌ای به شعاع $a=(3V/4\pi)^{1/3}\approx15.9$ نانومتر (قطر حدود ۳۲ نانومتر) است که دقیقاً همین مقدار طلا را در بر دارد. لازم به ذکر است که نسبت طول به عرض ۳ در اینجا برای استوانه‌ی کلاهک‌دار تعریف شده است و رابطه‌ی (\ref{eq:alpha-i}) فصل دوم، که برای بیضی‌گون نوشته شده، تنها به صورت کیفی برای آن معتبر است. محیط پس‌زمینه در تمامی شبیه‌سازی‌ها، خلأ/هوا با ضریب شکست واحد و ثابت ($n=1$) فرض شده است تا اثرات خالص نانوآنتن، بدون تداخل با پاشندگی محیط میزبان، مورد مطالعه قرار گیرد.

برای مدل‌سازی دقیق خواص اپتیکی نانومیله‌ی طلا، از داده‌های تجربی استاندارد جانسون و کریستی (\lr{Johnson and Christy}) استفاده شده است \cite{jc}. موتور شبیه‌سازی، تابع دی‌الکتریک مختلط طلا را در بازه‌ی طیفی شبیه‌سازی (۵۰۰ تا ۹۰۰ نانومتر) بر اساس الگوریتم معادلات دیفرانسیل کمکی (\lr{ADE}) و با استفاده از مدل چندضریبی بر روی این داده‌های تجربی برازش (\lr{Fit}) کرده است تا اثرات گذارهای بین‌نواری نوار $d$ به درستی در محاسبات زمانی لحاظ گردد \cite{maier,taflove}.

گسیلنده‌ی کوانتومی نیز با یک دوقطبی الکتریکی نقطه‌ای جایگزین گردیده است \cite{novotny}. در مدل‌سازی پایه (\lr{Reference Model})، به منظور رصد قوی‌ترین جفت‌شدگی ممکن با مد طولی، محور طولی نانومیله در راستای محور $z$ دستگاه مختصات شبیه‌سازی قرار دارد و دوقطبی با جهت‌گیری قطبشی موازی با همین محور ($z$)، در فاصله‌ی بسیار نزدیک ۳ نانومتری از رأس کلاهک بالایی و روی محور میله قرار داده شد. با این چیدمان، کل سامانه (میله و دوقطبی) حول محور $z$ تقارن استوانه‌ای دارد؛ بنابراین راستاهای $x$ و $y$ کاملاً هم‌ارزند. در سراسر این فصل، منظور از «گاف» ($d$) همین فاصله‌ی گسیلنده تا نزدیک‌ترین نقطه‌ی سطح فلز است. به منظور حفظ ثبات شرط کورانت و دقت محاسباتی در این فاصله‌ی کوچک، یک ناحیه‌ی شبکه‌بندی موضعی بسیار ظریف (\lr{Mesh Override}) با گام فضایی همسان‌گرد $\Delta x=\Delta y=\Delta z=1$ نانومتر در اطراف دوقطبی و کلاهک نانومیله تعریف گردید \cite{taflove}. این مش‌بندی متراکم، تضمین می‌کند که افت سریع (توانی) و تغییرات تند میدان نزدیک (وابستگی $1/d^3$ در تابع گرین) در فضای میان دوقطبی و سطح فلز، به صورت پیوسته و بدون پله‌پله شدن عددی (\lr{Staircase effect}) محاسبه شود \cite{novotny,taflove}.

\section{تحلیل طیفی فاکتور پورسل و تفکیک مدهای تشدید پلاسمونی}
مهم‌ترین معیار برای سنجش قدرت یک نانوآنتن در فشرده‌سازی میدان نوری و تغییر محیط الکترومغناطیسی پیرامون گسیلنده، بررسی طیف «فاکتور پورسل» ($F_P$) است \cite{purcell,novotny}. این کمیت بدون‌بُعد، نشان می‌دهد که حضور نانومیله تا چه حد توانسته است نرخ فروافت کل دوقطبی (مجموع کانال‌های تابشی و غیرتابشی) را نسبت به نرخ فروافت آن در فضای آزاد (خلأ) افزایش دهد \cite{novotny}. بر اساس قضایای الکترودینامیک، این فاکتور با محاسبه‌ی کل توان الکترومغناطیسی خارج‌شده از منبع ($P_{\sub{tot}}$) به دست می‌آید \cite{novotny}.

داده‌های استخراج‌شده از جعبه‌ی مانیتورینگ پورسل (محصورکننده‌ی منبع) برای گاف ۳ نانومتر، پرده از یک رفتار طیفی بسیار غنی و دوگانه برمی‌دارد. طیف محاسبه‌شده در بازه‌ی ۵۰۰ تا ۹۰۰ نانومتر (نمایش داده شده در شکل \ref{fig:purcell})، دو قله‌ی (\lr{Peak}) کاملاً متمایز با منشأ فیزیکی متفاوت را نشان می‌دهد \cite{maier,novotny}:

\begin{figure}[htbp]
\centering
\figbox{purcell_nm.png}{\textwidth}
\caption{طیف فاکتور پورسل ($F_P$) بر حسب طول موج ($\lambda$، از ۵۰۰ تا ۹۰۰ نانومتر) استخراج‌شده از داده‌های نرم‌افزار \lr{Lumerical FDTD} برای دوقطبی محوری در فاصله‌ی گاف ۳ نانومتری از نانومیله‌ی طلا. نمودار دو قله‌ی شاخص را در طول‌موج‌های ۵۱۵ نانومتر ($F_P\approx7200$، ناحیه‌ی جذب بین‌نواری و مدهای غیرتابشی) و ۶۰۷ نانومتر ($F_P\approx4211$، مد تشدید پلاسمون سطحی طولی) نشان می‌دهد.}
\label{fig:purcell}
\end{figure}

همان‌گونه که در شکل \ref{fig:purcell} مشاهده می‌گردد، محور افقی بیانگر طول موج (از ۵۰۰ تا ۹۰۰ نانومتر) و محور عمودی بیانگر مقدار بدون‌بُعد فاکتور پورسل است. تحلیل کمّی این منحنی، دو رژیم فیزیکی زیر را آشکار می‌سازد:

\begin{enumerate}
\item \textbf{قله‌ی بسامد بالا (جذب بین‌نواری و مدهای غیرتابشی):} قله‌ی نخست در طول‌موج حدود ۵۱۵ نانومتر (شکل \ref{fig:purcell}) ظاهر شده است و فاکتور پورسل را به حدود ۷۲۰۰ می‌رساند (یعنی دوقطبی ۷۲۰۰ بار سریع‌تر از فضای آزاد انرژی خود را تخلیه می‌کند) \cite{novotny}. این طول‌موج دقیقاً با آستانه‌ی گذارهای بین‌نواری طلا (حدود $2.4\,\mathrm{eV}$، معادل ۵۱۷ نانومتر) منطبق است؛ جایی که الکترون‌های نوار کاملاً پرشده‌ی $5d$ به حالت‌های خالی نوار $6sp$ برانگیخته می‌شوند و بخش موهومی تابع دی‌الکتریک طلا ($\epsilon_2$) به شدت افزایش می‌یابد \cite{maier,jc}. هرچند این ناحیه‌ی طیفی با محل تشدید عرضی نانومیله (و نیز تشدید نانوکره) هم‌پوشانی دارد، اما دوقطبی‌ای که روی محور میله و موازی با آن قرار دارد، بنا بر تقارن، قادر به تحریک مؤثر مد دوقطبی عرضی نیست \cite{maier}. بنابراین این قله‌ی عظیم، نه نشان‌دهنده‌ی تابش فوتون، بلکه ناشی از جذب مستقیم انرژی دوقطبی توسط گذارهای بین‌نواری و مدهای چندقطبی مرتبه‌بالای غیرتابشی در فاصله‌ی بسیار نزدیک است \cite{novotny,koenderink,ruppin}. در این طول‌موج، دوقطبی به جای تابش، به عنوان یک منبع حرارتی فوق‌سریع عمل کرده و انرژی خود را مستقیماً به شبکه‌ی فلزی تزریق می‌کند \cite{maier}.

\item \textbf{قله‌ی بسامد پایین (مد طولی و پنجره‌ی تابشی):} در مقابل، قله‌ی دوم که پس از یک کمینه‌ی موضعی در حدود ۵۸۰ نانومتر ($F_P\approx3100$) شکل گرفته است، در طول‌موج حدود ۶۰۷ نانومتر (ناحیه‌ی نارنجی طیف مرئی) ظاهر شده و فاکتور پورسلی در حدود ۴۲۰۰ (مقدار دقیق خوانده‌شده از نمودار: ۴۲۱۱) را به ثبت رسانده است (شکل \ref{fig:purcell}). این قله مستقیماً و منحصراً متناظر با «تشدید پلاسمون سطحی طولی» (\lr{Longitudinal LSPR}) در نانومیله است \cite{maier,novotny}. از نظر فیزیکی، در این طول‌موج، نوسانگرهای الکترونی در امتداد محور بلند نانومیله (محور ۶۰ نانومتری) به نوسان درمی‌آیند. از آنجا که فاصله‌ی قطب‌ها از یکدیگر بیشتر است، نیروی بازگرداننده‌ی (\lr{Restoring force}) کولنی که توسط بارهای مثبت پس‌زمینه به ابر الکترونی وارد می‌شود، ضعیف‌تر از حالت عرضی است \cite{maier}. همین نیروی بازگرداننده‌ی ضعیف‌تر است که باعث کاهش بسامد تشدید و انتقال قله به طول‌موج حدود ۶۰۷ نانومتر می‌شود \cite{maier}. در این ناحیه از طیف، گذارهای بین‌نواری طلا به شدت تضعیف شده و بخش موهومی تابع دی‌الکتریک ($\epsilon_2$) کاهش می‌یابد \cite{maier,jc}. رسیدن به فاکتور پورسل ۴۲۰۰ در این طول‌موج (و سپس افت تدریجی و مجانبی منحنی در طول‌موج‌های بالای ۷۰۰ نانومتر به دلیل دور شدن از شرط تشدید)، نشان می‌دهد که نانومیله‌ی طلا در فاصله‌ی ۳ نانومتری، میدان نزدیک را با شدتی بسیار زیاد (به واسطه‌ی اثر برق‌گیر در کلاهک) متمرکز کرده و یک مسیر بسیار سریع برای تخلیه‌ی انرژی دوقطبی فراهم آورده است؛ مسیری که این بار، پتانسیل بالایی برای هدایت تابشی دارد \cite{novotny,bharadwaj}.
\end{enumerate}

برای مقایسه، تشدید دوقطبی نانوکره‌ی طلای هم‌حجم در همین شرایط (خلأ، داده‌های جانسون-کریستی) بر اساس رابطه‌ی (\ref{eq:sphere-rates}) در حدود ۵۳۰ نانومتر قرار می‌گیرد؛ یعنی تنها در فاصله‌ی چند ده نانومتری از لبه‌ی گذارهای بین‌نواری و در همان ناحیه‌ی پراتلافی که در نانومیله تنها قله‌ی غیرتابشی را تولید کرده است \cite{bohren,kreibig}. به بیان دیگر، نانوکره امکان جدا کردن «پنجره‌ی تابشی» از «ناحیه‌ی جذب بین‌نواری» را ندارد، در حالی که نانومیله با انتقال تشدید طولی به حدود ۶۰۷ نانومتر، این دو را از هم جدا کرده است \cite{link,sonnichsen}.

\section{محاسبه‌ی توان تابشی از طریق مانیتورهای انتقال در فضای دوردست}
همان‌طور که آلبرت کوئندرینک نشان داده است، افزایش چشمگیر نرخ کل فروافت (فاکتور پورسل ۴۲۰۰) به هیچ وجه برای ارزیابی عملکرد یک نانوآنتن به عنوان «چشمه‌ی فوتون» کافی نیست \cite{koenderink}؛ بلکه باید مشخص شود این انرژی تخلیه‌شده چه مقدار از چنگال اتلاف اهمی طلا نجات داده و به صورت فوتون‌های مفید به فضای دوردست (\lr{Far-field}) تابش می‌کند \cite{novotny,koenderink}.

برای محاسبه‌ی کمّی این پارامتر حیاتی، یک «جعبه‌ی مانیتورینگ انتقال توان» (\lr{Transmission Box}) متشکل از ۶ مانیتور صفحه‌ای مجزا ($+x,-x,+y,-y,+z,-z$) در فضای دور از نانومیله و در مجاورت مرزهای جاذب \lr{PML} تعریف گردید \cite{novotny,taflove}. این صفحات، انتگرال مؤلفه‌ی عمودی بردار پوینتینگ ($\mathbf{S}$) را بر روی کل سطح خود محاسبه می‌کنند \cite{novotny}.

تحلیل داده‌های عددی استخراج‌شده از این ۶ مانیتور (نشان داده شده در شکل‌های \ref{fig:boxx}، \ref{fig:boxy} و \ref{fig:boxz})، رفتار تابشی سامانه‌ی جفت‌شده را به دقت روشن می‌سازد. در این نمودارها، علامت مثبت یا منفی صرفاً نشان‌دهنده‌ی جهت هندسی جریان توان نسبت به بردارهای یکه‌ی محورهای مختصات است.

\begin{figure}[htbp]
\centering
\begin{minipage}[t]{0.48\textwidth}\centering
\figbox{box_x1_nm.png}{\textwidth}\\[2pt]
{\small (آ) مانیتور وجه \lr{x1} (صفحه‌ی $-x$) با قله‌ی $T\approx-54.5$}
\end{minipage}\hfill
\begin{minipage}[t]{0.48\textwidth}\centering
\figbox{box_x2_nm.png}{\textwidth}\\[2pt]
{\small (ب) مانیتور وجه \lr{x2} (صفحه‌ی $+x$) با قله‌ی $T\approx+54.5$}
\end{minipage}
\caption{طیف انتقال توان تابشی ($T$) عبوری از وجوه جانبی عمود بر محور $x$ (\lr{box\_x1} و \lr{box\_x2}) بر حسب طول موج در نرم‌افزار لومریکال. محور میله و دوقطبی در راستای $z$ است.}
\label{fig:boxx}
\end{figure}

مطابق شکل \ref{fig:boxx}، در صفحات عمود بر محور $x$ (وجوه \lr{x1} و \lr{x2})، که از وجوه جانبی و موازی با محور میله هستند، طیف انتقال توان در طول‌موج حدود ۶۱۳ نانومتر، یعنی در ناحیه‌ی تشدید طولی، دارای یک قله‌ی بسیار تیز و تشدیدی است که در وجه \lr{x1} به مقدار $-54.5$ واحد و در وجه قرینه‌ی \lr{x2} به مقدار $+54.5$ واحد می‌رسد. نکته‌ی بسیار کلیدی در شکل \ref{fig:boxx} این است که در طول‌موج ۵۱۵ نانومتر، با وجود اینکه فاکتور پورسل در شکل \ref{fig:purcell} به ۷۲۰۰ رسیده بود، توان تابشی عبوری از مانیتورهای میدان دور تقریباً صفر (کمتر از ۱ واحد) است. این مقایسه به وضوح نشان می‌دهد که قله‌ی ۵۱۵ نانومتر یک قله‌ی کاملاً غیرتابشی و تاریک (\lr{Dark/Quenched}) است که تمام انرژی آن صرف گرم کردن نانومیله‌ی طلا می‌شود \cite{novotny,maier}.

\begin{figure}[htbp]
\centering
\begin{minipage}[t]{0.48\textwidth}\centering
\figbox{box_y1_nm.png}{\textwidth}\\[2pt]
{\small (آ) مانیتور وجه \lr{y1} (صفحه‌ی $-y$) با قله‌ی $T\approx-54.5$}
\end{minipage}\hfill
\begin{minipage}[t]{0.48\textwidth}\centering
\figbox{box_y2_nm.png}{\textwidth}\\[2pt]
{\small (ب) مانیتور وجه \lr{y2} (صفحه‌ی $+y$) با قله‌ی $T\approx+54.5$}
\end{minipage}
\caption{طیف انتقال توان تابشی ($T$) عبوری از وجوه جانبی راستای محور $y$ (\lr{box\_y1} و \lr{box\_y2}) بر حسب طول موج در نرم‌افزار لومریکال.}
\label{fig:boxy}
\end{figure}

بررسی مقادیر انتقال ($T$ - \lr{Transmission}) در راستای محور $y$ (شکل \ref{fig:boxy}) نیز نشان می‌دهد که توزیع فضایی توان به شدت متقارن است. در طول‌موج تشدید (حدود ۶۱۳ نانومتر)، توان عبوری از وجه $+y$ (\lr{box\_y2}) تا مقدار $+54.5$ واحد (نسبت به توان کل دوقطبی در خلأ) افزایش می‌یابد و به صورت کاملاً قرینه، وجه مقابل آن ($-y$ یا \lr{box\_y1})، توانی برابر با $-54.5$ را ثبت می‌کند که حاکی از خروج متقارن امواج در امتداد محور $y$ است. تساوی کامل منحنی‌های وجوه $x$ و $y$ (شکل‌های \ref{fig:boxx} و \ref{fig:boxy}) همان چیزی است که تقارن استوانه‌ای سامانه حول محور $z$ پیش‌بینی می‌کند و نشانه‌ی شکل‌گیری الگوی تابشی دوقطبی در صفحه‌ی عمود بر محور میله است \cite{novotny,maier}.

\begin{figure}[htbp]
\centering
\begin{minipage}[t]{0.48\textwidth}\centering
\figbox{box_z1_nm.png}{\textwidth}\\[2pt]
{\small (آ) مانیتور وجه \lr{z1} (صفحه‌ی $-z$) با کمینه‌ی $T\approx-17.5$}
\end{minipage}\hfill
\begin{minipage}[t]{0.48\textwidth}\centering
\figbox{box_z2_nm.png}{\textwidth}\\[2pt]
{\small (ب) مانیتور وجه \lr{z2} (صفحه‌ی $+z$) با بیشینه‌ی $T\approx+16.5$}
\end{minipage}
\caption{طیف انتقال توان تابشی ($T$) عبوری از وجوه تحتانی و فوقانی عمود بر محور $z$، یعنی وجوهی که در امتداد محور میله قرار دارند (\lr{box\_z1} و \lr{box\_z2})، بر حسب طول موج در نرم‌افزار لومریکال.}
\label{fig:boxz}
\end{figure}

همین تقارن ساختاری برای وجوه $z$ نیز در شکل \ref{fig:boxz} مشاهده می‌شود. در نمودار \lr{box\_z1} (وجه پایینی $-z$)، شار توان عبوری در طول موج حدود ۶۱۳ نانومتر به مقدار $-17.5$ واحد و در نمودار \lr{box\_z2} (وجه بالایی $+z$) به مقدار $+16.5$ واحد می‌رسد. سهم وجوه $z$ (حدود ۱۷ واحد) در مقایسه با وجوه جانبی $x$ و $y$ (هر کدام ۵۴/۵ واحد) به طور محسوسی کوچک‌تر است. این رفتار دقیقاً با هندسه‌ی سامانه سازگار است: محور میله و دوقطبی هر دو در راستای $z$ قرار دارند و مد طولی نانومیله مانند یک دوقطبی بزرگ در راستای $z$ تابش می‌کند. الگوی تابشی چنین دوقطبی‌ای متناسب با $\sin^2\theta$ است ($\theta$ زاویه نسبت به محور $z$)؛ یعنی در امتداد محور خود تقریباً تابش نمی‌کند و بیشینه‌ی تابش آن در صفحه‌ی عمود بر محور ($xy$) است \cite{novotny}. به همین دلیل، وجوه جانبی $x$ و $y$ سهم اصلی توان را دریافت می‌کنند و به دلیل تقارن استوانه‌ای حول محور $z$، سهم آن‌ها کاملاً برابر است (هر دو ۵۴/۵ واحد)؛ وجوه $z$ تنها پرتوهای مایل را دریافت می‌کنند و سهم کمتری دارند. اختلاف جزئی میان دو وجه $z$ ($17.5$ در برابر $16.5$) نیز طبیعی است، زیرا گسیلنده تنها در کنار کلاهک بالایی قرار دارد و سامانه نسبت به صفحه‌ی $z=0$ تقارن آینه‌ای ندارد.

با جمع‌بندی قدر مطلق انتگرال توان از تمامی این ۶ وجه
\[
|{-54.5}|+|{+54.5}|+|{-54.5}|+|{+54.5}|+|{-17.5}|+|{+16.5}|=252\,.\]
موتور محاسباتی \lr{FDTD} توان کل تابشی نانوآنتن ($P_{\sub{rad}}$) را استخراج می‌کند. نتایج نشان می‌دهد که طیف توان تابشی بر مد طولی منطبق است؛ با این تفاوت که قله‌ی آن (حدود ۶۱۳ نانومتر) چند نانومتر نسبت به قله‌ی فاکتور پورسل (۶۰۷ نانومتر) به سمت طول‌موج‌های بلندتر قرار دارد. این جابه‌جایی کوچک طبیعی است: فاکتور پورسل مجموع کانال‌های تابشی و غیرتابشی است و سهم غیرتابشی، به دلیل نزدیکی به لبه‌ی گذارهای بین‌نواری، با کاهش طول‌موج افزایش می‌یابد و قله‌ی نرخ کل را اندکی به سمت آبی می‌کشد، در حالی که قله‌ی توان تابشی تنها از مد طولی پیروی می‌کند \cite{novotny,koenderink}. مجموع مقادیر عبوری نشان می‌دهد که بیشینه‌ی توان تابشی سامانه در قله‌ی مد طولی، به عدد چشمگیر ۲۵۲ دست یافته است. مفهوم فیزیکی این عدد (۲۵۲) بسیار حائز اهمیت است: این عدد نشان می‌دهد که علی‌رغم حضور یک جسم فلزی اتلافی در همسایگی ۳ نانومتری دوقطبی، نانومیله نه تنها تمام انرژی را نبلعیده، بلکه به عنوان یک آنتن به شدت کارآمد عمل کرده و توانسته است نرخ تابش فوتون به فضای آزاد ($\Gamma_{\sub{rad}}$) را ۲۵۲ برابر زمانی که دوقطبی به تنهایی در خلأ قرار داشت، افزایش دهد \cite{novotny,anger}. این امر تأیید می‌کند که نانومیله دقیقاً در بسامد تشدید پلاسمونی طولی، نقش یک آنتن تابشی مؤثر را دارد و میدان نزدیک دوقطبی را با راندمان قابل‌توجهی به امواج انتشاریافته تبدیل می‌کند \cite{maier,novotny}. باید میان این دو کمیت تمایز گذاشت: تقویت تابشی ($\Gamma_{\sub{rad}}/\Gamma_0$) بیشینه‌ی آهنگ تولید فوتون را در شرایط اشباع تعیین می‌کند، در حالی که در یک چشمه‌ی تک‌فوتونی تحریک‌پالسی، که در هر پالس حداکثر یک فوتون تولید می‌شود، احتمال دریافت فوتون مفید در هر پالس برابر با بازده $\eta$ است \cite{koenderink,novotny}. به همین دلیل، در ادامه هر دو کمیت در کنار هم بررسی می‌شوند.

برای مقایسه، بر اساس رابطه‌ی (\ref{eq:sphere-rates})، نانوکره‌ی طلای هم‌حجم در همین فاصله (۳ نانومتر)، همین جهت‌گیری دوقطبی (عمود بر سطح) و همین طول‌موج (۶۰۷ نانومتر)، فاکتور پورسلی در حدود ۴۷۱ و تقویت تابشی‌ای تنها در حدود ۶/۹ برابر خلأ ایجاد می‌کند \cite{ruppin,mertens}. حتی در طول‌موج تشدید خود نانوکره (حدود ۵۳۰ نانومتر)، تقویت تابشی آن از حدود ۹ برابر فراتر نمی‌رود. بنابراین، با مقدار یکسان طلا، هندسه‌ی میله‌ای تقویت تابشی را حدود ۳۵ برابر افزایش داده است؛ نتیجه‌ای که با گزارش‌های تجربی مبنی بر تقویت فلورسانس تا حدود هزار برابر توسط تک‌نانومیله‌ها \cite{khatua} در برابر تقویت حداکثر چند ده برابری توسط نانوکره‌ها \cite{anger,kuhn} هم‌خوانی دارد.

\section{پویش فضایی گاف (\lr{Gap Sweep}) و رقابت میان تقویت تابشی و بازده کوانتومی}
یکی از چالش‌های بنیادین و شناخته‌شده در فیزیک نانواپتیک، وابستگی فوق‌العاده شدید برهم‌کنش‌های نور و ماده به فاصله‌ی فضایی میان گسیلنده و سطح فلز (طول گاف - $d$) است \cite{novotny,anger}. همان‌طور که در فصل سوم و در مبحث تابع گرین دایادیک اثبات شد، در رژیم میدان نزدیک، دامنه‌ی میدان پراکنده با ضریب معکوس مکعبی ($1/d^3$) تغییر می‌کند \cite{novotny}. این وابستگی نشان می‌دهد که جابه‌جایی‌های چند نانومتری می‌تواند دینامیک سامانه را کاملاً دگرگون سازد. به منظور یافتن نقطه‌ی کار بهینه‌ی چشمه‌های تک‌فوتونی، یک سری شبیه‌سازی‌های پارامتریک و سیستماتیک (\lr{Sweep}) برای گاف‌های ۳، ۵ و ۲۰ نانومتر انجام پذیرفت.

نتایج پویش فضایی نشان می‌دهد که با دور شدن دوقطبی از سطح کلاهک نانومیله، مقادیر مطلق تقویت (فاکتور پورسل کل) با شیب بسیار تندی افت می‌کنند \cite{novotny,maier}. در حالی که فاکتور پورسل در قله‌ی طولی برای گاف ۳ نانومتر برابر با ۴۲۰۰ بود، این مقدار با افزایش گاف به تنها ۵ نانومتر، به شدت سقوط کرده و به عدد ۱۴۹۶ می‌رسد. با ادامه‌ی این روند و قرارگیری گسیلنده در فاصله‌ی ۲۰ نانومتری، فاکتور پورسل به عدد ۹۱ تقلیل می‌یابد (البته باید توجه داشت که نرخ فروافت کل ۹۱ برابر خلأ، و نرخ تابشی حدود ۶۷ برابر خلأ، همچنان مقادیر قابل توجهی برای چشمه‌های فوتونی محسوب می‌شوند) \cite{novotny}. با وجود این تغییرات شدید دامنه‌ای، مکان طیفی قله‌ی تابشی پایداری بسیار بالایی از خود نشان داده و برای تمامی فواصل، تنها در حد چند نانومتر جابه‌جا می‌شود که نشان از استقلال بسامد \lr{LSPR} از بارگذاری محیطی دوقطبی دارد \cite{maier}.

اما نکته‌ی کلیدی، پنهان، و ارزشمند این پویش فضایی، در بررسی رفتار متناقض «بازده کوانتومی ظاهری» ($\eta$) نهفته است \cite{novotny,anger}. بازده کوانتومی از تقسیم توان تابشی به فضای دوردست (۲۵۲ در گاف ۳ نانومتر) بر توان کل خروجی از دوقطبی (۴۲۰۰ در گاف ۳ نانومتر) به دست می‌آید (رابطه‌ی (\ref{eq:eta})) \cite{novotny}. محاسبات نشان می‌دهد که در فواصل بسیار نزدیک (۳ تا ۵ نانومتر)، اگرچه فاکتور پورسل بسیار بزرگ است، اما بخش اعظم این انرژی (بیش از ۸۰ درصد) وارد کانال‌های اتلافی و مدهای پلاسمونی غیرمفیدی می‌شود که هرگز به فضا تابش نمی‌کنند \cite{novotny,maier}. در نتیجه، بازده تابشی ظاهری در گاف ۳ نانومتر تنها حدود ۶ درصد و در گاف ۵ نانومتر در حدود ۱۹ درصد است. این اعداد اثبات عددی وقوع پدیده‌ی «خاموشی غیرتابشی» (\lr{Quenching}) هستند؛ دوقطبی انرژی خود را هزاران بار سریع‌تر تخلیه می‌کند، اما در هر چرخه‌ی تحریک تنها حدود ۶ درصد احتمال تولید یک فوتون مفید وجود دارد و بقیه‌ی انرژی صرف گرم کردن ساختار می‌شود \cite{novotny,anger}.

با این وجود، هنگامی که دوقطبی به فواصل دورتر (گاف ۲۰ نانومتر) منتقل می‌شود، رفتار سامانه به طور اساسی تغییر می‌کند. تابع گرین پراکنده که نمایانگر جفت‌شدگی میدان نزدیک است، افت می‌کند، اما این افت برای کانال‌های غیرتابشی (تزریق به شبکه‌ی طلا) بسیار سریع‌تر از کانال‌های پراکندگی تابشی نانوآنتن اتفاق می‌افتد \cite{novotny,anger}. در نتیجه‌ی این عدم تقارن در نرخ افت، بازده کوانتومی سامانه‌ی جفت‌شده به طور چشمگیری، از ۱۹ درصد (در ۵ نانومتر) به حدود ۷۴ درصد (در گاف ۲۰ نانومتر) می‌رسد \cite{novotny}. در این فاصله، تقویت نرخ تابشی ($\eta F_P$) همچنان حدود ۶۷ برابر خلأ است. خلاصه‌ای از مقادیر عددی استخراج‌شده در این پویش فضایی در جدول \ref{tab:gap} گردآوری شده است.

\begin{table}[htbp]
\centering
\caption{خلاصه‌ی نتایج عددی استخراج‌شده از پویش فضایی گاف (در مد طولی تشدید، $\lambda\approx607$ نانومتر).}
\label{tab:gap}
\renewcommand{\arraystretch}{1.3}
\resizebox{\textwidth}{!}{%
\begin{tabular}{cccc}
\toprule
فاصله‌ی گاف $d$ (نانومتر) & فاکتور پورسل کل ($F_P$) & بازده کوانتومی ظاهری ($\eta$) & رژیم غالب فیزیکی \\
\midrule
۳ & ۴۲۰۰ & حدود ۶ درصد & خاموشی غیرتابشی شدید (\lr{Quenching}) \\
۵ & ۱۴۹۶ & حدود ۱۹ درصد & غلبه‌ی اتلاف اهمی بر تابش \\
۲۰ & ۹۱ & حدود ۷۴ درصد & ناحیه‌ی بهینه‌ی کارکرد (رژیم طلایی آنتن) \\
\bottomrule
\end{tabular}}
\end{table}

این نتیجه‌ی تحلیلی، مهم‌ترین دستاورد فیزیکی این پژوهش است. این تحلیل به طور قاطع نشان می‌دهد که برای مهندسی و ساخت عملی چشمه‌های تک‌فوتونیِ پلاسمونی، چسباندن گسیلنده به سطح فلز (جهت رسیدن به پورسل‌های هزار برابری) یک خطای مهندسی است \cite{novotny,koenderink}. در عوض، \textbf{گاف‌های فضایی میانی (حدود ۱۰ تا ۲۰ نانومتر)} به عنوان «محدوده‌ی طلایی» طراحی معرفی می‌گردند \cite{novotny,anger}. نقطه‌ی محاسبه‌شده‌ی این محدوده، گاف ۲۰ نانومتر است؛ قرار گرفتن فواصل حدود ۱۰ نانومتر در همین ناحیه، از روند پیوسته‌ی تغییرات میان گاف‌های ۵ و ۲۰ نانومتر و نیز از رفتار مشابه نانوکره (جدول \ref{tab:sphere}) انتظار می‌رود و محاسبه‌ی مستقیم آن پیشنهاد می‌شود. در این محدوده‌ی فضایی، یک مصالحه و تعادل فیزیکی مطلوب برقرار می‌شود: سامانه‌ی نانوآنتن همچنان قادر است نرخ گسیل (سرعت چشمه) را ده‌ها برابر فضای آزاد افزایش دهد (کاهش طول عمر به مقیاس پیکوثانیه)، و در عین حال، در انتهای این بازه (گاف ۲۰ نانومتر) با برخورداری از بازده کوانتومی حدود ۷۴ درصد، تضمین می‌کند که از هر ۱۰ چرخه‌ی تحریک گسیلنده، حدود ۷ فوتون با موفقیت به فضای دوردست تابش شده و توسط اپتیک جمع‌آوری قابل دریافت باشند \cite{novotny,maier}.

برای سنجش نقش هندسه در این مصالحه، همین پویش برای نانوکره‌ی طلای هم‌حجم (قطر حدود ۳۲ نانومتر) در شرایط کاملاً یکسان (خلأ، داده‌های جانسون-کریستی، دوقطبی عمود بر سطح و همان فواصل) با استفاده از حل تحلیلی دقیق شبه‌استاتیک (رابطه‌ی (\ref{eq:sphere-rates})، شامل تمام مرتبه‌های چندقطبی) محاسبه شد \cite{ruppin,gersten,mertens}. نتایج در جدول \ref{tab:sphere} در کنار نتایج نانومیله آمده است. مقایسه در دو طول‌موج انجام شده است: طول‌موج کاری گسیلنده (۶۰۷ نانومتر، تشدید نانومیله) و طول‌موج تشدید خود نانوکره (حدود ۵۳۰ نانومتر)، تا هر ساختار در شرایط مطلوب خود نیز سنجیده شود.

\begin{table}[htbp]
\centering
\caption{مقایسه‌ی نانومیله (شبیه‌سازی \lr{FDTD}) با نانوکره‌ی طلای هم‌حجم (حل تحلیلی دقیق شبه‌استاتیک) در شرایط یکسان. $\Gamma_{\sub{rad}}/\Gamma_0=\eta F_P$ تقویت نرخ تابشی نسبت به خلأ است.}
\label{tab:sphere}
\renewcommand{\arraystretch}{1.3}
\resizebox{\textwidth}{!}{%
\begin{tabular}{c|ccc|ccc|cc}
\toprule
 & \multicolumn{3}{c|}{نانومیله، حدود ۶۰۷ نانومتر} & \multicolumn{3}{c|}{نانوکره، ۶۰۷ نانومتر} & \multicolumn{2}{c}{نانوکره، ۵۳۰ نانومتر} \\
$d$ (نانومتر) & $F_P$ & $\Gamma_{\sub{rad}}/\Gamma_0$ & $\eta$ (درصد) & $F_P$ & $\Gamma_{\sub{rad}}/\Gamma_0$ & $\eta$ (درصد) & $\Gamma_{\sub{rad}}/\Gamma_0$ & $\eta$ (درصد) \\
\midrule
۳  & ۴۲۰۰ & ۲۵۲ & ۶  & ۴۷۱ & ۶/۹ & ۱/۵ & ۹/۰ & ۰/۴ \\
۵  & ۱۴۹۶ & ۲۸۴ & ۱۹ & ۱۰۴ & ۴/۹ & ۴/۷ & ۶/۱ & ۱/۲ \\
۲۰ & ۹۱ & ۶۷ & ۷۴ & ۲/۰ & ۱/۵ & ۷۷ & ۱/۷ & ۲۵ \\
\bottomrule
\end{tabular}}
\end{table}

این مقایسه سه نتیجه‌ی روشن دارد. نخست، روند کلی در هر دو هندسه یکسان است: خاموشی شدید در فواصل زیر ۵ نانومتر و افزایش بازده با دور شدن گسیلنده؛ این همان رفتار جهانی ناشی از جمله‌ی $1/d^3$ تابع گرین است که در آزمایش‌های کلاسیک روی نانوکره‌ها نیز مشاهده شده است \cite{anger,kuhn}. دوم، در هر فاصله‌ی یکسان، تقویت تابشی نانومیله حدود ۳۵ تا ۶۰ برابر نانوکره است (برای مثال ۲۵۲ در برابر ۶/۹ در گاف ۳ نانومتر و ۶۷ در برابر ۱/۵ در گاف ۲۰ نانومتر)؛ و حتی اگر نانوکره در طول‌موج تشدید خود به کار رود، تقویت تابشی آن از حدود ۹ برابر فراتر نمی‌رود و بازده آن به دلیل هم‌پوشانی با گذارهای بین‌نواری بسیار کم می‌شود \cite{sonnichsen,link}. سوم، در گاف ۲۰ نانومتر بازده نانوکره (۷۷ درصد) کمی بیشتر از نانومیله (۷۴ درصد) است، اما این بازده بالا تنها به این دلیل است که نانوکره در این فاصله تقریباً هیچ اثری بر گسیلنده ندارد (تقویت تابشی ۱/۵ برابر)؛ بنابراین بازده بالا به تنهایی معیار مناسبی نیست و باید همراه با تقویت تابشی گزارش شود \cite{koenderink}. به این ترتیب، «محدوده‌ی طلایی» فواصل میانی، که در آن هم بازده بالا و هم تقویت تابشی چشمگیر هم‌زمان حاصل می‌شوند، ویژگی هندسه‌ی میله‌ای است و برای نانوکره‌ی هم‌حجم وجود ندارد.

\section{بررسی ناهمسانگردی جهتی (\lr{Directional Anisotropy}) در پاسخ نانوآنتن}
یکی از ویژگی‌های بنیادینی که نانومیله‌های طلا را از نانوکره‌های متقارن متمایز می‌سازد، حساسیت فوق‌العاده شدید پاسخ اپتیکی آن‌ها به زاویه و قطبش منبع تحریک‌کننده است \cite{maier,novotny}. برای ارزیابی کمّی این ناهمسانگردی در رژیم جفت‌شدگی میدان نزدیک، شبیه‌سازی‌ها با تغییر جهت‌گیری بردار گشتاور دوقطبی (از حالت موازی با محور طولی نانومیله، به حالت عرضی عمود بر محور نانومیله) تکرار گردید.

بررسی نتایج در حالت تحریک عرضی (\lr{Transverse dipole orientation}) نشان داد که در طول‌موج تشدید طولی (حدود ۶۰۷ نانومتر)، فاکتور پورسل از ۴۲۰۰ (حالت طولی) به حدود ۱۰۲ کاهش می‌یابد ($4200/102\approx41$) و تقویت توان تابشی به کمتر از یک می‌رسد؛ یعنی تقریباً تمام همین ۱۰۲ برابر افزایش نرخ فروافت، از کانال غیرتابشی (جذب در فلز نزدیک گسیلنده) عبور می‌کند و تابش مفید حتی از خلأ نیز ضعیف‌تر است \cite{novotny,maier}. در این حالت، دوقطبی عملاً توانایی تحریک کردن بارهای الکترونی در امتداد محور بلند نانومیله را ندارد و تشدید طولی کاملاً خاموش (\lr{Dark}) باقی می‌ماند \cite{maier}.

با استخراج داده‌ها و محاسبه‌ی نسبت فاکتور پورسل در حالت تحریک طولی به حالت تحریک عرضی در بسامد تشدید (حدود ۶۰۷ نانومتر)، عدد چشمگیر ۴۱ به دست آمد. به منظور درک بهتر این عدد، همین نسبت برای یک نانوکره‌ی هم‌حجم شبیه‌سازی و محاسبه شد. نتایج گروه کنترل (نانوکره) نشان داد که به دلیل تقارن کروی، نسبت پاسخ‌ها در تمام بازه‌ی طیفی تقریباً حول عدد ۲/۴ ثابت است و وابستگی قطبشی ضعیفی دارد (که آن هم صرفاً ناشی از حضور دوقطبی در یک سمت کره است) \cite{novotny,maier}. این مقدار با حل تحلیلی دقیق (رابطه‌ی (\ref{eq:sphere-rates})) نیز به خوبی تأیید می‌شود: نسبت فاکتور پورسل دوقطبی شعاعی به دوقطبی مماسی برای نانوکره‌ی هم‌حجم در گاف ۳ نانومتر و طول‌موج ۶۰۷ نانومتر برابر با ۲/۴۱ به دست می‌آید \cite{ruppin}. این نسبت از ترکیب مدهای چندقطبی با وزن‌های $(n+1)^2$ و $n(n+1)/2$ ناشی می‌شود و مقداری میان ۲ و ۴ دارد؛ یعنی منشأ آن صرفاً جهت‌گیری نسبت به سطح است و نه وجود یک مد ممتاز.

مقایسه‌ی نسبت ۴۱ در نانومیله با نسبت ۲/۴ در نانوکره، به طور قطعی ثابت می‌کند که ناهمسانگردی شدید مشاهده‌شده در نانومیله، منحصراً ناشی از تحریک یک مد پلاسمونی قدرتمند در امتداد محور ممتاز ساختار است و صرفاً یک اثر هندسی ساده‌ی نزدیک‌سطحی محسوب نمی‌شود \cite{maier,giannini}. این ویژگی برجسته نشان می‌دهد که نانومیله‌ی پلاسمونی طلا نه تنها نرخ گسیل را افزایش می‌دهد، بلکه به عنوان یک «فیلتر قطبشی نانومتری» عمل کرده و تنها فوتون‌هایی را که دارای قطبش هم‌راستا با محور آنتن هستند، با راندمان بالا به فضای دوردست پرتاب می‌کند؛ ویژگی‌ای که برای کدگذاری اطلاعات در مخابرات کوانتومی بر اساس قطبش فوتون‌ها، بسیار ایده‌آل است \cite{novotny,maier}.

\section{تحلیل حساسیت، همگرایی مش و پایداری محاسبات (\lr{Mesh Convergence})}
در تمامی پژوهش‌های مبتنی بر الکترودینامیک محاسباتی، اعتبارسنجی نتایج و اطمینان از استقلال داده‌های فیزیکی از پارامترهای گسسته‌سازی (روش‌های مش‌بندی)، امری کاملاً ضروری است \cite{taflove}. در روش \lr{FDTD}، به خصوص در مسائلی که با پدیده‌ی میدان نزدیک و مرزهای منحنی (مانند کلاهک نیم‌کروی نانومیله) سروکار دارند، خطاهای موسوم به «پلکانی شدن مرزها» (\lr{Staircase error}) و عدم تقارن شبکه می‌توانند نتایج را مخدوش سازند \cite{taflove}.

به منظور سنجش اعتبار و استواری محاسبات انجام‌شده در این فصل، یک آزمون سیستماتیک «همگرایی شبکه» (\lr{Mesh Convergence Test}) طراحی و اجرا گردید \cite{taflove}. در این آزمون، شبیه‌سازی حالت پایه (گاف ۳ نانومتر) به طور متوالی با استفاده از گام‌های شبکه‌بندی فضایی درشت تا بسیار ظریف، شامل مقادیر $\Delta x=3$، ۲، ۱/۵ و ۱ نانومتر اجرا شد.

تحلیل طیف‌های به دست آمده از این آزمون چهارمرحله‌ای حاکی از آن است که مؤلفه‌های فیزیکی اساسی سامانه، پایداری بی‌نظیری از خود نشان می‌دهند. مکان طیفی تشدیدها کاملاً قفل شده و با تغییر اندازه‌ی مش، تنها در یک حاشیه‌ی خطای بسیار محدود ($607\pm4$ نانومتر) نوسان می‌کند. به طور مشابه، بازده تابشی محاسبه‌شده در بسامد تشدید نیز در تمامی شبکه‌بندی‌ها، استواری خود را حفظ کرده و در بازه‌ی حدود ۶ تا ۹ درصد باقی می‌ماند.

با این حال، باید به یک واقعیت ذاتی در شبیه‌سازی‌های \lr{FDTD} منابع نقطه‌ای در فواصل فوق‌کوچک اذعان نمود. به دلیل پدیده‌ای که در متون حل‌کننده‌های عددی به عنوان «قرارگیری زیرسلولی چشمه» (\lr{Sub-cell source placement}) شناخته می‌شود، و همچنین حساسیت شدید (توانی) میدان نزدیک به تغییرات فضایی در مقیاس‌های کمتر از طول شبکه، ارتفاع مطلق قله‌های فاکتور پورسل در گاف ۳ نانومتر دارای یک عدم قطعیت سیستماتیک عددی در حدود ۴۰ درصد میان شبکه‌ی ۳ نانومتری و شبکه‌ی ۱ نانومتری است (برای مثال، قله‌ی ۵۱۵ نانومتر از مقدار حدود ۵۰۰۰ به ۷۲۰۰ رشد می‌کند) \cite{taflove,novotny}.

کاهش بیشتر ابعاد شبکه (برای مثال از ۱ به ۰/۵ نانومتر) هزینه‌ی محاسباتی را افزایش می‌دهد، اما این رشد نمایی نیست. بر اساس شرط کورانت (رابطه‌ی \lr{CFL})، گام زمانی با اندازه‌ی شبکه رابطه‌ی خطی دارد ($\Delta t\propto\Delta x$)؛ بنابراین نصف شدن اندازه‌ی شبکه در هر سه بُعد، تعداد سلول‌ها را ۸ برابر و تعداد گام‌های زمانی را ۲ برابر می‌کند و حجم کل محاسبات حداکثر ۱۶ برابر می‌شود (رشدی از مرتبه‌ی $\Delta x^{-4}$) \cite{taflove}. در عمل، چون شبکه‌ی ریز تنها در ناحیه‌ی کوچک \lr{Mesh Override} اطراف دوقطبی و کلاهک اعمال شده است و نه در کل فضای شبیه‌سازی، افزایش زمان اجرا کمتر از این حد (چند برابر) خواهد بود؛ هرچند نیاز به حافظه و زمان همچنان قابل توجه است. افزون بر این، حتی با شبکه‌ی ریزتر، خطای ناشی از قرارگیری زیرسلولی چشمه در فاصله‌ی ۳ نانومتری به طور کامل حذف نمی‌شود. به همین دلیل، مقادیر مطلق قله‌های پورسل در کوچک‌ترین گاف (۷۲۰۰ و ۴۲۰۰) باید با عدم قطعیتی در حد چند ده درصد تفسیر شوند \cite{taflove}. از این رو، استدلال‌های فیزیکی، مهندسی و نتیجه‌گیری‌های ارائه‌شده در سرتاسر این پژوهش، کاملاً آگاهانه تدوین شده و تماماً بر پایه‌ی متغیرهای پایدار و همگرا (مانند مکان طیفی تشدیدها، روند تغییرات در پویش فضایی، بازده‌های تابشی نسبی و نسبت‌های جهتی) بنا شده‌اند تا از خطاهای سیستماتیک ناشی از گسسته‌سازی فضا مصون بمانند.

حل تحلیلی نانوکره نیز به عنوان یک آزمون مستقل از گسسته‌سازی به کار گرفته شد. نانوکره، برخلاف نانومیله، حل دقیق دارد و بنابراین هیچ خطای مش در آن وجود ندارد. تطابق نسبت جهتی شبیه‌سازی‌شده برای نانوکره (۲/۴) با مقدار تحلیلی (۲/۴۱) نشان می‌دهد که تنظیمات مدل \lr{FDTD} (مدل ماده، منبع دوقطبی و جعبه‌های مانیتورینگ) درست پیاده‌سازی شده‌اند و تفاوت‌های مشاهده‌شده میان دو هندسه، فیزیکی‌اند و نه عددی \cite{ruppin,mertens}.

% ======================================== فصل ۶ ======================================================
\chapter{نتیجه‌گیری و پیشنهادات}

\section{نتیجه‌گیری کلی}
این پایان‌نامه با هدف بررسی، تحلیل و بهینه‌سازی برهم‌کنش میان یک گسیلنده‌ی کوانتومی منفرد و یک نانوآنتن پلاسمونی (نانومیله‌ی طلا) تدوین و اجرا گردید \cite{novotny,maier}. چالش بنیادین در طراحی و توسعه‌ی چشمه‌های تک‌فوتونی کارآمد مبتنی بر پلاسمونیک، غلبه بر پدیده‌ی مخرب «خاموشی غیرتابشی» (\lr{Non-radiative Quenching}) در فواصل فوق‌کوچک نانومتری از سطح فلز است؛ ناحیه‌ای که در آن، اتلاف اهمی طلا (جذب درون‌نواری الکترون‌های آزاد و، در طول‌موج‌های کوتاه‌تر، گذارهای بین‌نواری) بر تقویت تابشی غلبه کرده و انرژی گسیلنده را به حرارت تبدیل می‌کنند \cite{novotny,anger}.

برای کالبدشکافی دقیق این پدیده، از یک مدل الکترودینامیکی سه‌بعدی بدون تقریب‌های تحلیلی و مبتنی بر روش تفاضل محدود در حوزه‌ی زمان (\lr{3D-FDTD}) استفاده شد \cite{taflove}. با جایگزینی گسیلنده با یک نوسانگر دوقطبی کلاسیک و بهره‌گیری از نظریه‌ی تابع گرین دایادیک و قانون طلایی فرمی، سامانه‌ی جفت‌شده در هر دو حوزه‌ی زمان و بسامد به صورت جامع مورد کاوش قرار گرفت \cite{novotny,maier}. برای آنکه سهم هندسه‌ی میله‌ای از سهم خود ماده‌ی طلا تفکیک شود، تمامی نتایج کلیدی با یک نانوکره‌ی طلای هم‌حجم (قطر حدود ۳۲ نانومتر)، که حل تحلیلی دقیق دارد و مرجع کلاسیک این حوزه است، در شرایط یکسان مقایسه شد \cite{ruppin,gersten,mertens}.

نتایج حاصل از این شبیه‌سازی‌های پیشرفته به روشنی نشان داد که رویکرد سنتی مبنی بر «نزدیک کردن هرچه بیشتر گسیلنده به سطح فلز برای رسیدن به بالاترین فاکتور پورسل»، یک استراتژی ناکارآمد برای ساخت چشمه‌های نوری است \cite{koenderink,novotny}. در عوض، این پژوهش نشان داد که با مهندسی دقیق فاصله (فضای گاف) میان گسیلنده و نانوآنتن، می‌توان به یک مصالحه‌ی فیزیکی مطلوب میان «سرعت تخلیه‌ی انرژی» و «بازده کوانتومی» دست یافت. مقایسه با نانوکره نیز نشان داد که این مصالحه‌ی مطلوب، که در آن تقویت تابشی بزرگ و بازده بالا هم‌زمان حاصل می‌شوند، ویژگی هندسه‌ی میله‌ای است: با مقدار یکسان طلا، نانوکره در هیچ فاصله‌ای در بازه‌ی ۳ تا ۲۰ نانومتر و در هیچ طول‌موجی به تقویت تابشی بیش از حدود ۹ برابر نمی‌رسد، در حالی که نانومیله ده‌ها تا صدها برابر تقویت ایجاد می‌کند \cite{novotny,maier,khatua}. بنابراین عملکرد نانومیله‌ی طلا نه تنها به عنوان یک تشدیدگر موضعی، بلکه به عنوان یک هدایتگر فوتونی فوق‌سریع و کارآمد قابل اتکاست.

\section{دستاوردهای کلیدی پژوهش}
مهم‌ترین نتایج، یافته‌های عددی و استنتاج‌های فیزیکی استخراج‌شده از این شبیه‌سازی‌ها را می‌توان در چهار محور اصلی زیر طبقه‌بندی و خلاصه نمود:

\begin{enumerate}
\item \textbf{شناسایی و تفکیک دقیق مدهای تشدیدی پلاسمونی:} تحلیل طیفی سامانه‌ی جفت‌شده نشان داد که نانومیله‌ی طلا (به طول ۶۰ و قطر ۲۰ نانومتر) در پاسخ الکترومغناطیسی دو قله‌ی کاملاً متمایز دارد \cite{maier}. قله‌ی نخست، در حوالی ۵۱۵ نانومتر و بر لبه‌ی گذارهای بین‌نواری طلا (نوار $d$ به $sp$) قرار دارد؛ اگرچه فاکتور پورسل در این قله در گاف ۳ نانومتر به حدود ۷۲۰۰ می‌رسد، اما ماهیت آن کاملاً اتلافی (تاریک) است و توان تابشی متناظر با آن تقریباً صفر است \cite{maier,novotny}. در مقابل، مد تشدیدی طولی که در طول‌موج حدود ۶۰۷ نانومتر ظاهر می‌گردد، با ایجاد فاکتور پورسل ۴۲۰۰، بستر اصلی و قدرتمندی را برای هدایت فوتون‌ها به فضای آزاد (با ثبت بیشینه‌ی توان تابشی ۲۵۲ برابر خلأ) فراهم می‌آورد. نانوکره‌ی هم‌حجم، در مقابل، تنها یک تشدید دوقطبی یگانه در حدود ۵۳۰ نانومتر دارد که دقیقاً بر همان ناحیه‌ی پراتلاف بین‌نواری منطبق است؛ به همین دلیل «پنجره‌ی تابشی» و «ناحیه‌ی جذب» در نانوکره از یکدیگر جداشدنی نیستند و طول‌موج تشدید آن نیز عملاً قابل تنظیم نیست \cite{kreibig,link,sonnichsen}. بنابراین شکستن تقارن کروی و انتقال تشدید طولی به طول‌موج‌های بلندتر، که با تغییر نسبت طول به عرض به سادگی ممکن است، شرط اول طراحی نانوآنتن کارآمد است \cite{novotny,maier}.

\item \textbf{یافتن فاصله‌ی بهینه (نقطه‌ی کار ایده‌آل برای چشمه‌های تک‌فوتونی):} پویش سیستماتیک فضای گاف (\lr{Gap Sweep}) نشان داد که قرارگیری گسیلنده در فواصل کمتر از ۵ نانومتر، به دلیل رشد معکوس مکعبی تابع گرین پراکنده ($1/d^3$) در رژیم میدان نزدیک، منجر به سلطه‌ی مطلق اتلاف فلز بر سامانه شده و بازده تابشی ظاهری را به ۶ درصد (در گاف ۳ نانومتر) و حداکثر ۱۹ درصد (در گاف ۵ نانومتر) تقلیل می‌دهد \cite{novotny,anger}. دستاورد بنیادین و مهندسی این پژوهش، نشان دادن این موضوع بود که فواصل فضایی میانی (حدود ۱۰ تا ۲۰ نانومتر، با نقطه‌ی محاسبه‌شده در گاف ۲۰ نانومتر)، نقطه‌ی طلایی طراحی هستند. در این محدوده، اگرچه فاکتور پورسل مطلق از ۱۴۹۶ (در گاف ۵ نانومتر) به ۹۱ (در گاف ۲۰ نانومتر) کاهش می‌یابد، اما سامانه از چنگال اتلاف اهمی رها شده و بازده کوانتومی نانوآنتن تا حدود ۷۴ درصد (در گاف ۲۰ نانومتر) افزایش می‌یابد، در حالی که نرخ تابشی همچنان حدود ۶۷ برابر خلأ است \cite{novotny}. نانوکره‌ی هم‌حجم، که بر اساس حل تحلیلی دقیق در همین فواصل محاسبه شد، همان روند کلی (خاموشی در فواصل کوتاه و افزایش بازده با فاصله) را دارد که نشان می‌دهد این روند ناشی از قانون عمومی $1/d^3$ است و به هندسه وابسته نیست \cite{ruppin,anger}. اما مقدار تقویت تابشی در هر فاصله‌ی یکسان، در نانومیله حدود ۳۵ تا ۶۰ برابر نانوکره است؛ برای مثال، در گاف ۲۰ نانومتر نانوکره با وجود بازده ۷۷ درصد، نرخ تابشی را تنها ۱/۵ برابر تقویت می‌کند. به عبارت دیگر، «محدوده‌ی طلایی» فواصل میانی، که در انتهای آن (گاف ۲۰ نانومتر) بیش از ۷۰ درصد از انرژی دوقطبی با سرعتی ده‌ها برابر خلأ به شکل فوتون‌های قابل جمع‌آوری تابش می‌شود، در مقایسه‌ی هم‌حجم تنها برای نانومیله وجود دارد \cite{novotny,koenderink}.

\item \textbf{کمّی‌سازی ناهمسانگردی جهتی شدید (\lr{Anisotropy}):} محاسبه و مقایسه‌ی سیستماتیک پاسخ ساختار به جهت‌گیری‌های گوناگون گسیلنده نشان داد که نانومیله رفتاری به شدت وابسته به جهت دوقطبی دارد \cite{maier,novotny}. نسبت فاکتور پورسل در حالت تحریک طولی به حالت تحریک عرضی، عدد ۴۱ را به ثبت رساند. در مقایسه با نانوذرات کروی هم‌حجم (که این نسبت در آن‌ها حدود ۲/۴ است، مقداری که با حل تحلیلی دقیق ۲/۴۱ نیز تأیید شد)، این عدد نشان‌دهنده‌ی شکست کامل تقارن کروی است \cite{maier,ruppin}. این یافته نشان می‌دهد که نانومیله می‌تواند به عنوان یک «فیلتر قطبشی» عمل کرده و فوتون‌ها را با قطبشی کاملاً مشخص و جهت‌دار ساطع کند که برای رمزنگاری کوانتومی بر پایه‌ی قطبش، یک ویژگی بی‌نظیر محسوب می‌شود \cite{novotny}.

\item \textbf{پایداری محاسباتی و تفکیک خطاهای عددی:} بررسی‌های دقیق همگرایی شبکه (\lr{Mesh convergence}) نشان داد که در شبیه‌سازی‌های \lr{FDTD}، ردیابی فواصل فوق‌کوچک همواره با عدم قطعیت‌هایی در ارتفاع مطلق قله‌های پورسل (به دلیل پدیده‌ی قرارگیری زیرسلولی چشمه) همراه است \cite{taflove}. با این وجود، این پژوهش نشان داد که مؤلفه‌های فیزیکی حیاتی برای طراحی مدارهای فوتونیک، مانند «مکان طیفی تشدیدها»، «بازده‌های تابشی» و «نسبت‌های جهتی»، از ثبات و پایداری عددی فوق‌العاده بالایی برخوردارند و استدلال‌های ارائه‌شده در این پایان‌نامه، که بر همین کمیت‌ها بنا شده‌اند، عمدتاً مستقل از خطاهای گسسته‌سازی شبکه هستند. افزون بر این، نانوکره به دلیل داشتن حل تحلیلی دقیق، به عنوان یک آزمون مستقل از مش به کار رفت: تطابق نسبت جهتی شبیه‌سازی‌شده‌ی نانوکره (۲/۴) با مقدار تحلیلی (۲/۴۱)، صحت پیاده‌سازی مدل ماده، منبع دوقطبی و جعبه‌های مانیتورینگ را تأیید می‌کند و نشان می‌دهد که تفاوت‌های گزارش‌شده میان نانومیله و نانوکره، منشأ فیزیکی دارند و ناشی از خطای عددی نیستند \cite{ruppin,mertens}.
\end{enumerate}

\section{پیشنهادات برای پژوهش‌های آتی}
با توجه به نتایج ارزشمند به دست آمده در این پایان‌نامه و وسعت افق‌های موجود در حوزه‌ی الکترودینامیک کوانتومی درون‌کاواکی (\lr{CQED}) و نانوفوتونیک، مسیرهای پژوهشی زیر برای گسترش، تعمیق و تکمیل این مطالعات به محققان و دانشجویان آینده پیشنهاد می‌گردد:

\begin{itemize}
\item \textbf{بررسی هندسه‌های جفت‌شده و آنتن‌های دایمر (\lr{Dimers}):} پژوهش حاضر بر روی یک نانومیله‌ی منفرد (\lr{Monomer}) متمرکز بود. پیشنهاد می‌شود در کارهای آینده، ساختار دوتایی نانومیله‌ها (به شکل نانوآنتن‌های پاپیونی یا \lr{Bowtie}) شبیه‌سازی گردد \cite{maier,giannini}. قرار دادن گسیلنده دقیقاً در مرکز گاف میان دو نانومیله می‌تواند به دلیل بروز اثر «هیبریدی‌شدن پلاسمون‌ها» (\lr{Plasmon Hybridization})، فشرده‌سازی میدان و در نتیجه فاکتور پورسل را تا چندین مرتبه‌ی بزرگی فراتر از نتایج این پایان‌نامه افزایش دهد \cite{maier,novotny}.

\item \textbf{تغییر جنس نانوآنتن و کاوش در طیف‌های فرابنفش:} طلا به دلیل پایداری شیمیایی بالا گزینه‌ای عالی است، اما اتلاف ناشی از گذارهای بین‌نواری ($5d\to6sp$) آن در طول‌موج‌های کمتر از حدود ۵۵۰ نانومتر بسیار شدید است \cite{maier,jc}. تکرار این شبیه‌سازی‌ها و کدهای محاسباتی با استفاده از فلزاتی مانند نقره (\lr{Ag}) و آلومینیوم (\lr{Al}) قویاً پیشنهاد می‌شود. نقره به دلیل داشتن اتلاف اهمی بسیار کمتر در کل ناحیه‌ی مرئی، و آلومینیوم برای کاربرد در ناحیه‌ی فرابنفش (\lr{UV})، می‌توانند بازده‌ی تابشی سامانه و طیف عملکردی چشمه‌های تک‌فوتونی را به طول‌موج‌های کوتاه‌تر ارتقا دهند \cite{maier}.

\item \textbf{ورود به رژیم جفت‌شدگی قوی (\lr{Strong Coupling Regime}):} در این پژوهش، سامانه تماماً در «رژیم جفت‌شدگی ضعیف» بررسی شد؛ جایی که نانوآنتن صرفاً نرخ فروافت برگشت‌ناپذیر دوقطبی را از طریق اثر پورسل تسریع می‌کند \cite{novotny}. پیشنهاد می‌شود با طراحی ساختارهایی با حجم مدی حتی کوچک‌تر (مانند نانوگاف‌های پیکومتری) و استفاده از گسیلنده‌هایی با گشتاور دوقطبی بسیار بزرگ، ورود سامانه به رژیم جفت‌شدگی قوی مدل‌سازی شود. در این رژیم، انرژی به صورت همدوس (\lr{Coherent}) میان گسیلنده و پلاسمون مبادله شده و پدیده‌های کوانتومی جذابی مانند «شکافتگی رابی خلأ» (\lr{Vacuum Rabi Splitting}) در طیف انتقال سامانه ظهور خواهند کرد \cite{novotny,maier}.

\item \textbf{اثرات چندگسیلنده‌ای (\lr{Multi-emitter effects}) و همدوسی فضایی:} قرار دادن هم‌زمان چندین دوقطبی کوانتومی در مجاورت نانومیله و بررسی برهم‌کنش آن‌ها از طریق مدهای پلاسمونی مشترک، یک افق تحقیقاتی بسیار جذاب است \cite{novotny}. این شبیه‌سازی‌ها می‌توانند پدیده‌های کوانتومی جمعی، مانند «ابرتابش دیک» (\lr{Dicke Superradiance}) و تأثیر میدان‌های پلاسمونی مقید بر قفل‌شدگی فاز (\lr{Phase-locking}) و همدوسی فضایی مجموعه‌ای از گسیلنده‌ها را مورد کاوش قرار دهند \cite{novotny,maier}.
\end{itemize}

% ======================================== منابع ======================================================
\chapter*{منابع و مآخذ}
\addcontentsline{toc}{chapter}{منابع و مآخذ}
\markboth{منابع و مآخذ}{منابع و مآخذ}
\begin{latin}
\renewcommand{\bibname}{\rl{منابع و مآخذ}}
\renewcommand{\chapter}[2]{}   % عنوان فارسی بالا چاپ شده است؛ عنوان لاتین تکراری حذف می‌شود
\begin{thebibliography}{99}
\bibitem{novotny} L. Novotny and B. Hecht, \textit{Principles of Nano-Optics}, 2nd ed. Cambridge, UK: Cambridge University Press, 2012.
\bibitem{maier} S. A. Maier, \textit{Plasmonics: Fundamentals and Applications}. New York, NY, USA: Springer, 2007.
\bibitem{purcell} E. M. Purcell, ``Spontaneous emission probabilities at radio frequencies,'' \textit{Physical Review}, vol. 69, no. 11--12, p. 681, 1946.
\bibitem{drexhage} K. H. Drexhage, ``Influence of a dielectric interface on fluorescence decay time,'' \textit{Journal of Luminescence}, vol. 1--2, pp. 693--701, 1970.
\bibitem{yablonovitch} E. Yablonovitch, ``Inhibited spontaneous emission in solid-state physics and electronics,'' \textit{Physical Review Letters}, vol. 58, no. 20, p. 2059, 1987.
\bibitem{john} S. John, ``Strong localization of photons in certain disordered dielectric superlattices,'' \textit{Physical Review Letters}, vol. 58, no. 23, p. 2486, 1987.
\bibitem{barnes} W. L. Barnes, A. Dereux, and T. W. Ebbesen, ``Surface plasmon subwavelength optics,'' \textit{Nature}, vol. 424, no. 6950, pp. 824--830, 2003.
\bibitem{anger} P. Anger, P. Bharadwaj, and L. Novotny, ``Enhancement and quenching of single-molecule fluorescence,'' \textit{Physical Review Letters}, vol. 96, no. 11, p. 113002, 2006.
\bibitem{bharadwaj} P. Bharadwaj, B. Deutsch, and L. Novotny, ``Optical antennas,'' \textit{Advances in Optics and Photonics}, vol. 1, no. 3, pp. 438--483, 2009.
\bibitem{rogobete} L. Rogobete, F. Kaminski, M. Agio, and V. Sandoghdar, ``Design of plasmonic nanoantennae for enhancing spontaneous emission,'' \textit{Optics Letters}, vol. 32, no. 12, pp. 1623--1625, 2007.
\bibitem{koenderink} A. F. Koenderink, ``On the use of Purcell factors for plasmon antennas,'' \textit{Optics Letters}, vol. 35, no. 24, pp. 4208--4210, 2010.
\bibitem{giannini} V. Giannini, A. I. Fern\'andez-Dom\'inguez, S. C. Heck, and S. A. Maier, ``Plasmonic nanoantennas: fundamentals and their use in controlling the radiative properties of nanoemitters,'' \textit{Chemical Reviews}, vol. 111, no. 6, pp. 3888--3912, 2011.
\bibitem{pelton} M. Pelton, ``Modified spontaneous emission in nanophotonic structures,'' \textit{Nature Photonics}, vol. 9, no. 7, pp. 427--435, 2015.
\bibitem{jc} P. B. Johnson and R. W. Christy, ``Optical constants of the noble metals,'' \textit{Physical Review B}, vol. 6, no. 12, pp. 4370--4379, 1972.
\bibitem{yee} K. S. Yee, ``Numerical solution of initial boundary value problems involving Maxwell's equations in isotropic media,'' \textit{IEEE Transactions on Antennas and Propagation}, vol. 14, no. 3, pp. 302--307, 1966.
\bibitem{taflove} A. Taflove and S. C. Hagness, \textit{Computational Electrodynamics: The Finite-Difference Time-Domain Method}, 3rd ed. Norwood, MA, USA: Artech House, 2005.
\bibitem{berenger} J.-P. Berenger, ``A perfectly matched layer for the absorption of electromagnetic waves,'' \textit{Journal of Computational Physics}, vol. 114, no. 2, pp. 185--200, 1994.
\bibitem{mie} G. Mie, ``Beitr\"age zur Optik tr\"uber Medien, speziell kolloidaler Metall\"osungen,'' \textit{Annalen der Physik}, vol. 330, no. 3, pp. 377--445, 1908.
\bibitem{bohren} C. F. Bohren and D. R. Huffman, \textit{Absorption and Scattering of Light by Small Particles}. New York, NY, USA: Wiley, 1983.
\bibitem{kreibig} U. Kreibig and M. Vollmer, \textit{Optical Properties of Metal Clusters}. Berlin, Germany: Springer, 1995.
\bibitem{link} S. Link and M. A. El-Sayed, ``Spectral properties and relaxation dynamics of surface plasmon electronic oscillations in gold and silver nanodots and nanorods,'' \textit{Journal of Physical Chemistry B}, vol. 103, no. 40, pp. 8410--8426, 1999.
\bibitem{sonnichsen} C. S\"onnichsen, T. Franzl, T. Wilk, G. von Plessen, J. Feldmann, O. Wilson, and P. Mulvaney, ``Drastic reduction of plasmon damping in gold nanorods,'' \textit{Physical Review Letters}, vol. 88, no. 7, p. 077402, 2002.
\bibitem{gersten} J. Gersten and A. Nitzan, ``Spectroscopic properties of molecules interacting with small dielectric particles,'' \textit{Journal of Chemical Physics}, vol. 75, no. 3, pp. 1139--1152, 1981.
\bibitem{ruppin} R. Ruppin, ``Decay of an excited molecule near a small metal sphere,'' \textit{Journal of Chemical Physics}, vol. 76, no. 4, pp. 1681--1684, 1982.
\bibitem{kuhn} S. K\"uhn, U. H\aa kanson, L. Rogobete, and V. Sandoghdar, ``Enhancement of single-molecule fluorescence using a gold nanoparticle as an optical nanoantenna,'' \textit{Physical Review Letters}, vol. 97, no. 1, p. 017402, 2006.
\bibitem{mertens} H. Mertens, A. F. Koenderink, and A. Polman, ``Plasmon-enhanced luminescence near noble-metal nanospheres: Comparison of exact theory and an improved Gersten and Nitzan model,'' \textit{Physical Review B}, vol. 76, no. 11, p. 115123, 2007.
\bibitem{khatua} S. Khatua, P. M. R. Paulo, H. Yuan, A. Gupta, P. Zijlstra, and M. Orrit, ``Resonant plasmonic enhancement of single-molecule fluorescence by individual gold nanorods,'' \textit{ACS Nano}, vol. 8, no. 5, pp. 4440--4449, 2014.
\end{thebibliography}
\end{latin}

\end{document}
